% Arthur Cayley, Collected Mathematical Papers, Vol. II, Pages 261--312
% Transcribed from: collectedmathema02cayluoft.pdf
% Papers: [141] A Second Memoir Upon Quantics (cont.)
%         [142] Numerical Tables Supplementary to Second Memoir on Quantics
%         [143] Tables of the Covariants M to W of the Binary Quintic
%         [144] A Third Memoir Upon Quantics

\documentclass[11pt,a4paper]{article}
\usepackage{amsmath}
\usepackage[utf8]{inputenc}
\usepackage[T1]{fontenc}
\usepackage{lmodern}
\usepackage{amsmath,amssymb,amsthm}
\usepackage{array,longtable,booktabs,multirow}
\usepackage{geometry}
\geometry{margin=1in}
\usepackage{fancyhdr}
\usepackage{graphicx}
\usepackage{enumitem}
\usepackage{titlesec}

% Custom commands for Cayley's notation
\newcommand{\ind}{\operatorname{ind.}}
\newcommand{\tab}[1]{\ \ \ \ \ \ \text{#1}}
\newcommand{\caybetween}{\mathbin{\raisebox{0.5ex}{\scriptsize$\pitchfork$}}}
\renewcommand{\between}{\caybetween}

% Page style
\pagestyle{fancy}
\fancyhf{}
\renewcommand{\headrulewidth}{0pt}
\fancyhead[C]{\thepage}

% Reduce paragraph indent
\setlength{\parindent}{1em}
\setlength{\parskip}{0.5ex}

\begin{document}
\setcounter{page}{261}

% =====================================================================
% PAPER [141] --- A SECOND MEMOIR UPON QUANTICS (continued from p. 260)
% =====================================================================

\begin{center}
\textsc{A Second Memoir Upon Quantics.}
\end{center}

\noindent and from this, by subtracting from each coefficient the coefficient which immediately precedes it, we form the table:

\begin{center}
\begin{tabular}{*{16}{c}}
\multicolumn{15}{r}{1} & (0)\\[4pt]
\multicolumn{13}{r}{1} & \multicolumn{2}{c}{0}\phantom{)} & (1)\\[4pt]
\multicolumn{11}{r}{1} & \multicolumn{3}{c}{0\quad 1\quad 0} & (2)\\[4pt]
\multicolumn{9}{r}{1} & \multicolumn{5}{c}{0\quad 1\quad 1\quad 1\quad 0} & (3)\\[4pt]
\multicolumn{7}{r}{1} & \multicolumn{7}{c}{0\quad 1\quad 1\quad 1\quad 1\quad 0\quad 1} & (4)\\[4pt]
\multicolumn{5}{r}{1} & \multicolumn{9}{c}{0\quad 1\quad 1\quad 1\quad 1\quad 1\quad 1\quad 0}\\[4pt]
\multicolumn{3}{r}{1} & \multicolumn{11}{c}{0\quad 1\quad 1\quad 1\quad 1\quad 2\quad 1\quad 1\quad 0} & (6)\\
\end{tabular}
\end{center}

The successive lines fix the number and character of the covariants of the degrees $0$, $1$, $2$, $3$, \&c. The line $(0)$, if this were to be interpreted, would show that there is a single covariant of the degree $0$; this covariant is of course merely the absolute constant unity, and may be excluded. The line $(1)$ shows that there is a single covariant of the degree $1$, viz.\ a covariant of the order $3$; this is the cubic itself, which I represent by $U$. The line $(2)$ shows that there are two asyzygetic covariants of the degree $2$, viz.\ one of the order $6$, this is merely $U^{2}$, and one of the order $2$, this I represent by $H$. The line $(3)$ shows that there are three asyzygetic covariants of the degree $3$, viz.\ one of the order $9$, this is $U^{3}$; one of the order $5$, this is $UH$, and one of the order $3$, this I represent by $\Phi$. The line $(4)$ shows that there are five asyzygetic covariants of the degree $4$, viz.\ one of the order $12$, this is $U^{4}$; one of the order $8$, this is $U^{2}H$; one of the order $6$, this is $H^{2}$; and one of the order $0$, i.e.\ an invariant, this I represent by $\nabla$. The line $(5)$ shows that there are six asyzygetic covariants of the degree $5$, viz.\ one of the order $15$, this is $U^{5}$; one of the order $11$, this is $U^{3}H$; one of the order $9$, this is $U^{2}\Phi$; one of the order $7$, this is $UH^{2}$; one of the order $5$, this is $H\Phi$; and one of the order $3$, this is $\nabla\,U$. The line $(6)$ shows that there are $8$ asyzygetic covariants of the degree $6$, viz.\ one of the order $18$, this is $U^{6}$; one of the

\clearpage
\setcounter{page}{262}

order $14$, this is $U^{4}H$; one of the order $12$, this is $U^{2}\Phi$; one of the order $10$, this is $U^{2}H^{2}$; one of the order $8$, this is $UH\Phi$; two of the order $6$ (i.e.\ the three covariants $H^{2}$, $\Phi^{2}$ and $\nabla\,U^{2}$ are not asyzygetic, but are connected by a single linear equation or syzygy), and one of the order $2$, this is $\nabla H$. We are thus led to the irreducible covariants $U$, $H$, $\Phi$, $\nabla$ connected by a linear equation or syzygy between $H^{2}$, $\Phi^{2}$ and $\nabla\,U^{2}$, and this is in fact the complete system of irreducible covariants; $\nabla$ is therefore the only invariant.

\textbf{36.} The asyzygetic covariants are of the form $U^{p}H^{q}\nabla^{r}$, or else of the form $U^{p}H^{q}\nabla^{r}\Phi$; and since $U$, $H$, $\nabla$ are of the degrees $1$, $2$, $4$ respectively, and $\Phi$ is of the degree $3$, the number of asyzygetic covariants of the degree $m$ of the first form is equal to the coefficient of $x^{m}$ in $1 \div (1-x)\,(1-x^{2})\,(1-x^{4})$, and the number of the asyzygetic covariants of the degree $m$ of the second form is equal to the coefficient of $x^{m}$ in $x^{3} \div (1-x)\,(1-x^{2})\,(1-x^{4})$. Hence the total number of asyzygetic covariants is equal to the coefficient of $x^{m}$ in $(1+x^{3}) \div (1-x)\,(1-x^{2})\,(1-x^{4})$, or what is the same thing, in
\[
\frac{1-x^{6}}{(1-x)\,(1-x^{2})^{2}\,(1-x^{4})}\,;
\]
and conversely, if this expression for the number of the asyzygetic covariants of the degree $m$ were established independently, it would follow that the irreducible invariants were four in number, and of the degrees $1$, $2$, $3$, $4$ respectively, but connected by an equation of the degree $6$. As regards the invariants, every invariant is of the form $\nabla^{p}$, i.e.\ the number of asyzygetic invariants of the degree $m$ is equal to the coefficient of $x^{m}$ in $\dfrac{1}{1-x^{4}}$, and conversely, from this expression it would follow that there was a single irreducible invariant of the degree $4$.

\textbf{37.} In the case of a quartic, the function to be developed is:
\[
\frac{1}{(1-z)\,(1-xz)\,(1-x^{2}z)\,(1-x^{3}z)\,(1-x^{4}z)}\,,
\]
and the coefficients are given by the table:

\begin{center}
\begin{tabular}{*{17}{c}}
\multicolumn{16}{r}{1} & (0)\\[4pt]
\multicolumn{14}{r}{1} & \multicolumn{2}{c}{1\quad 1} & (1)\\[4pt]
\multicolumn{12}{r}{1} & \multicolumn{4}{c}{1\quad 2\quad 2\quad 3} & (2)\\[4pt]
\multicolumn{10}{r}{1} & \multicolumn{6}{c}{1\quad 2\quad 3\quad 4\quad 5\quad 6} & (3)\\[4pt]
\multicolumn{8}{r}{1} & \multicolumn{8}{c}{1\quad 2\quad 3\quad 5\quad 6\quad 8\quad 9}\\[4pt]
\multicolumn{6}{r}{1} & \multicolumn{10}{c}{1\quad 2\quad 3\quad 5\quad 7\quad 9\quad 11\quad 14} & (5)\\[4pt]
\multicolumn{4}{r}{1} & \multicolumn{12}{c}{1\quad 2\quad 3\quad 5\quad 6\quad 9\quad 11\quad 14\quad 16\quad 19\quad 20} & (6)\\
\end{tabular}
\end{center}

\clearpage
\setcounter{page}{263}

and subtracting from each coefficient the coefficient immediately preceding it, we have the table:

\begin{center}
\begin{tabular}{*{17}{c}}
\multicolumn{16}{r}{1} & (0)\\[4pt]
\multicolumn{14}{r}{1} & \multicolumn{2}{c}{0\quad 0} & (1)\\[4pt]
\multicolumn{12}{r}{1} & \multicolumn{4}{c}{0\quad 1\quad 0\quad 1} & (2)\\[4pt]
\multicolumn{10}{r}{1} & \multicolumn{6}{c}{0\quad 1\quad 1\quad 1\quad 1\quad 0} & (3)\\[4pt]
\multicolumn{8}{r}{1} & \multicolumn{8}{c}{0\quad 1\quad 1\quad 2\quad 1\quad 2\quad 0\quad 1} & (4)\\[4pt]
\multicolumn{6}{r}{1} & \multicolumn{10}{c}{0\quad 1\quad 1\quad 2\quad 2\quad 2\quad 2\quad 0}\\[4pt]
\multicolumn{4}{r}{1} & \multicolumn{12}{c}{0\quad 1\quad 1\quad 2\quad 1\quad 3\quad 1\quad 3\quad 1\quad 2\quad 0\quad 2} & (6)\\
\end{tabular}
\end{center}

the examination of which will show that we have for the quartic the following irreducible covariants, viz.\ the quartic itself $U$; an invariant of the degree $2$, which I represent by $I$; a covariant of the order $4$ and of the degree $2$, which I represent by $H$; an invariant of the degree $3$, which I represent by $J$; and a covariant of the order $6$ and the degree $3$, which I represent by $\Phi$; but that the irreducible covariants are connected by an equation of the degree $6$, viz.\ there is a linear equation or syzygy between $\Phi^{2}$, $I^{3}H^{3}$, $IJH^{2}U$, $IJ^{2}HU^{2}$ and $J^{3}U^{3}$; this is in fact the complete system of the irreducible covariants of the quartic: the only irreducible invariants are the invariants $I$, $J$.

\textbf{38.} The asyzygetic covariants are of the form $U^{p}I^{q}H^{r}J^{s}$, or else of the form $U^{p}I^{q}H^{r}J^{s}\Phi$, and the number of the asyzygetic covariants of the degree $m$ is equal to the coefficient of $x^{m}$ in $(1+x^{3}) \div (1-x)\,(1-x^{2})^{2}\,(1-x^{3})$, or what is the same thing, in
\[
\frac{1-x^{6}}{(1-x)\,(1-x^{2})^{2}\,(1-x^{3})^{2}}\,,
\]
and the asyzygetic invariants are of the form $I^{p}J^{q}$, and the number of the asyzygetic invariants of the degree $m$ is equal to the coefficient of $x^{m}$ in $1 \div (1-x^{2})\,(1-x^{3})$. Conversely, if these formul\ae\ were established, the preceding results as to the form of the system of the irreducible covariants or of the irreducible invariants, would follow.

\textbf{39.} In the case of a quintic, the function to be developed is
\[
\frac{1}{(1-z)\,(1-xz)\,(1-x^{2}z)\,(1-x^{3}z)\,(1-x^{4}z)\,(1-x^{5}z)}\,;
\]
and the coefficients are given by the table:

\begin{center}
\begin{tabular}{*{19}{c}}
\multicolumn{18}{r}{1} & (0)\\[4pt]
\multicolumn{16}{r}{1} & \multicolumn{2}{c}{1\quad 1} & (1)\\[4pt]
\multicolumn{14}{r}{1} & \multicolumn{4}{c}{1\quad 2\quad 2\quad 3\quad 3} & (2)\\[4pt]
\multicolumn{12}{r}{1} & \multicolumn{6}{c}{1\quad 2\quad 3\quad 4\quad 5\quad 6\quad 6} & (3)\\[4pt]
\multicolumn{10}{r}{1} & \multicolumn{8}{c}{1\quad 2\quad 3\quad 5\quad 6\quad 8\quad 9\quad 11\quad 11\quad 12} & (4)\\[4pt]
\multicolumn{8}{r}{1} & \multicolumn{10}{c}{1\quad 2\quad 3\quad 5\quad 7\quad 9\quad 11\quad 14\quad 16\quad 18\quad 19\quad 20} & (5)\\
\end{tabular}
\end{center}

\clearpage
\setcounter{page}{264}

and subtracting from each coefficient the one which immediately precedes it, we have the table:

\begin{center}
\begin{tabular}{*{19}{c}}
\multicolumn{18}{r}{1} & (0)\\[4pt]
\multicolumn{16}{r}{1} & \multicolumn{2}{c}{0\quad 0} & (1)\\[4pt]
\multicolumn{14}{r}{1} & \multicolumn{4}{c}{0\quad 1\quad 0\quad 1\quad 0} & (2)\\[4pt]
\multicolumn{12}{r}{1} & \multicolumn{6}{c}{0\quad 1\quad 1\quad 1\quad 1\quad 1\quad 0} & (3)\\[4pt]
\multicolumn{10}{r}{1} & \multicolumn{8}{c}{0\quad 1\quad 1\quad 2\quad 1\quad 2\quad 1\quad 2\quad 0\quad 1} & (4)\\[4pt]
\multicolumn{8}{r}{1} & \multicolumn{10}{c}{0\quad 1\quad 1\quad 2\quad 2\quad 2\quad 2\quad 3\quad 2\quad 2\quad 1\quad 1} & (5)\\
\end{tabular}
\end{center}

We thus obtain the following irreducible covariants, viz.:

Of the degree $1$; a single covariant of the order $5$, this is the quintic itself.

Of the degree $2$; two covariants, viz.\ one of the order $6$, and one of the order $2$.

Of the degree $3$; three covariants, viz.\ one of the order $9$, one of the order $5$, and one of the order $3$.

Of the degree $4$; three covariants, viz.\ one of the order $6$, one of the order $4$, and one of the order $0$ (an invariant).

Of the degree $5$; three covariants, viz.\ one of the order $7$, one of the order $3$, and one of the order $1$ (a linear covariant).

These covariants are connected by a single syzygy of the degree $5$ and of the order $11$; in fact, the table shows that there are only two asyzygetic covariants of this degree and order; but we may, with the above-mentioned irreducible covariants of the degrees, $1$, $2$, $3$ and $4$, form three covariants of the degree $5$ and the order $11$; there is therefore a syzygy of this degree and order.

\clearpage
\setcounter{page}{265}

\textbf{40.} It is proper to remark, that the property that the roots of the equation $J=0$ are all real, does not in an obvious manner lead to any conclusion as to the reality of the roots of the resolvent equation. In fact, suppose that the roots of the equation are $\alpha$, $\beta$, $\gamma$, $\delta$, and that $I$ and $J$ are the two invariants of the biquadratic function $(\alpha,\beta,\gamma,\delta\,\raisebox{0.5ex}{$\pitchfork$}\, x, y)^{4}$, so that $I=0$ is the condition in order that there may be two equal roots; then if $H=(\alpha,\beta,\gamma,\delta\,\raisebox{0.5ex}{$\pitchfork$}\, x, y)^{4}$, the equation $H=0$ may be considered as a biquadratic equation the roots of which are the anharmonic ratios of the four roots $\alpha$, $\beta$, $\gamma$, $\delta$; and the invariants of this biquadratic function are in fact $I^{2}$ and $I^{3}-27J^{2}$; whence the equation $J=0$ gives $I^{3}-27J^{2}=I^{3}$, so that the condition $J=0$ expresses that the biquadratic equation $H=0$ has equal roots. And since the roots of $H=0$ are real, it follows that the equation in the anharmonic ratios has equal roots; but this does not seem to lead to any conclusion as to the reality of the roots of the resolvent equation.

\textbf{41.} A biquadratic equation $U=(a,b,c,d,e\,\raisebox{0.5ex}{$\pitchfork$}\, x,y)^{4}=0$ has been said to be in the canonical form when the second and fourth coefficients vanish, i.e.\ when $b=0$, $d=0$, and it has been shown that this can in general be done. The reduction depends on the solution of a cubic equation, and to effect it we have only to write $x+\lambda y$ for $x$; then if $a'$ is properly determined, the second coefficient will vanish identically; and if $\lambda$ is properly determined, the fourth coefficient will vanish; moreover the process shows that the equation is expressible in the canonical form in three different ways, corresponding to the three roots of the cubic equation. And the cubic equation for the determination of $a'$ is the resolvent above referred to; but I have not succeeded in making the proof depend upon the form of the equation of the anharmonic ratios.

\textbf{42.} The form of the biquadratic equation may be exhibited in a very simple manner by means of the invariant $J$; in fact, we have $I=0$ for the biquadratic function $(\alpha,\beta,\gamma,\delta\,\raisebox{0.5ex}{$\pitchfork$}\, x, y)^{4}$, and this gives the before-mentioned condition in order that the four roots may be equal two and two. But for the discriminant, we have $I^{3}-27J^{2}=0$; and if $I=0$, $J=0$, the four roots are all equal; hence $J=0$, if $I$ is not $=0$, is the condition in order that the roots may be equal in pairs.

\clearpage
\setcounter{page}{266}

\textbf{43.} The preceding investigations show the manner in which the general theory of the covariants of a binary quantic gives rise to the special results which have been established in relation to binary quantics of particular orders. I will consider, in the first place, the general question of the linear transformation of a binary quantic. This part of the subject is of course identical with the theory of the invariants of a binary quantic; and if the quantic is represented by
\[
U=(a,b,c,\ldots\,\raisebox{0.5ex}{$\pitchfork$}\, x,y)^{n},
\]
and if the new variables $x'$, $y'$ are given by
\[
x=\alpha x'+\beta y',\qquad y=\gamma x'+\delta y',
\]
then if $U$ is transformed into
\[
U'=(a',b',c',\ldots\,\raisebox{0.5ex}{$\pitchfork$}\, x',y')^{n},
\]
we have between the coefficients $a$, $b$, $c$,\ldots\ of the original function and the coefficients $a'$, $b'$, $c'$,\ldots\ of the transformed function, a system of $n+1$ equations, which may be written
\begin{align*}
a' &= (\alpha\,\raisebox{0.5ex}{$\pitchfork$}\, a,b,c,\ldots)^{n},\\
b' &= (\alpha\,\raisebox{0.5ex}{$\pitchfork$}\, a,b,c,\ldots)^{n-1}(\gamma\,\raisebox{0.5ex}{$\pitchfork$}\, a,b,c,\ldots),\\
c' &= (\alpha\,\raisebox{0.5ex}{$\pitchfork$}\, a,b,c,\ldots)^{n-2}(\gamma\,\raisebox{0.5ex}{$\pitchfork$}\, a,b,c,\ldots)^{2},\\
&\ \,\vdots
\end{align*}
and from these equations we may eliminate the quantities $\alpha$, $\beta$, $\gamma$, $\delta$, and thus obtain the invariants of the function; but the process is in general a very complicated one.

\textbf{44.} The question may be simplified as follows: writing $x'$, $y'$ in place of $x$, $y$, and considering $a'$, $b'$, $c'$,\ldots\ as given functions of $\alpha$, $\beta$, $\gamma$, $\delta$, we have
\[
(a,b,c,\ldots\,\raisebox{0.5ex}{$\pitchfork$}\, \alpha x'+\beta y',\, \gamma x'+\delta y')^{n}
 = (a',b',c',\ldots\,\raisebox{0.5ex}{$\pitchfork$}\, x',y')^{n},
\]
and this being so, if we operate with
\[
\Omega'=\frac{d}{dx'}\,\frac{d}{dy}-\frac{d}{dy'}\,\frac{d}{dx},
\]
the operation on the right-hand side gives zero; and attending to the formula
\[
\frac{d}{dx}=\alpha\,\frac{d}{dx'}+\gamma\,\frac{d}{dy'},\qquad
\frac{d}{dy}=\beta\,\frac{d}{dx'}+\delta\,\frac{d}{dy'},
\]
the operation on the left-hand side becomes
\[
(\alpha\delta-\beta\gamma)
\left(\frac{d}{dx'}\,\frac{d}{dy'}-\frac{d}{dy'}\,\frac{d}{dx'}\right)
\cdot(a,b,c,\ldots\,\raisebox{0.5ex}{$\pitchfork$}\, x',y')^{n},
\]
which vanishes identically; hence the equation is satisfied identically, and we thus see that $a'$, $b'$, $c'$,\ldots\ are such that the equation in $(x',y')$ is satisfied identically.

\clearpage
\setcounter{page}{267}

\textbf{45.} I use the term ``invariant'' to denote a function of the coefficients which possesses the property of invariance under linear transformation; but it is to be observed that every such function is not necessarily an invariant in the strict sense of the word. A function of the coefficients which is equal to the same function of the transformed coefficients multiplied by a power of the modulus (or determinant of transformation), is an invariant; but there may be functions which possess a different property, being equal to the same function of the transformed coefficients multiplied by a power of the modulus and also by a factor depending on the coefficients of the substitution. Such functions might be termed ``seminvariants''; but I shall not at present enter upon the theory of these functions.

\textbf{46.} I remark that if an invariant of a binary quantic is expressed as a function of the differences of the roots, then the invariant remains unaltered when any linear substitution is performed upon the variables. Suppose that the roots of the equation $U=0$ are $\alpha_{1}$, $\alpha_{2}$,\ldots, $\alpha_{n}$, then if we write $x=\alpha_{1}y$,\ldots, we may express any coefficient as a symmetric function of the differences of the roots; and any invariant expressed as a function of these differences will remain unaltered by any linear transformation of the variables. This is in fact the distinctive property of invariants.

\textbf{47.} I will now return to the general theory of covariants. Considering the quantic $U$, and forming with it the system of covariants, we have a number of covariantive equations, each of which equated to zero gives rise to a corresponding covariantive curve; and the system of these curves constitutes what may be termed the ``covariantive pencil'' of the quantic. The theory of the covariantive pencil presents many interesting questions, which I do not however propose to discuss in the present Memoir.

\textbf{48.} I will conclude with a few remarks on the expression of the covariants by means of the notation of bipartite forms. This notation is as follows: if $\alpha$, $\beta$ are symbols of substitution, and $(\alpha\,\raisebox{0.5ex}{$\pitchfork$}\, x,y)^{n}$, $(\beta\,\raisebox{0.5ex}{$\pitchfork$}\, x,y)^{n}$ are the same quantic, then the operator
\[
(\alpha\beta'-\alpha'\beta)^{p}(\alpha\,\raisebox{0.5ex}{$\pitchfork$}\, x,y)^{n-p}(\beta\,\raisebox{0.5ex}{$\pitchfork$}\, x,y)^{n-p}
\]
is a covariant of the quantic. In particular, if $p=1$, we have the canonizant; if $p=2$, we have a covariant of the order $2n-4$; and so on. The theory of these forms presents many interesting particular cases, which I reserve for future investigation.

\clearpage
\setcounter{page}{268}

\textbf{49.} The theory of the binary quantic leads naturally to the theory of ternary quantics and quantics in general; but the investigation of the covariants and invariants of ternary quantics presents difficulties of a much more formidable character than those which occur in the case of binary quantics. The leading coefficients of a ternary quantic are not, as in the case of a binary quantic, connected by a single differential equation, but by a system of partial differential equations; and the theory of the integration of these equations is very imperfectly understood. I hope however to be able to resume the consideration of this subject in a future Memoir.

\textbf{50.} I collect here for convenience of reference the principal formul\ae\ relating to the binary quantic:
\begin{alignat*}{2}
&(a,b,c,\ldots\,\raisebox{0.5ex}{$\pitchfork$}\, x,y)^{n}\qquad &&\text{(the quantic itself)},\\
&(ab)^{2}&&(=\text{Hessian}),\\
&(ab)^{2}(ac)(bd)(cd)^{2}&&(=\text{canonizant of deg.\ 3}),\\
&(ab)^{2}(ac)^{2}(bd)(cd)^{3}&&(=\text{canonizant of deg.\ 4}),\\
&\multicolumn{2}{l}{\&c.\qquad \&c.}
\end{alignat*}
where the symbols $(ab)$, $(ac)$,\ldots\ denote the determinants $\alpha\beta'-\alpha'\beta$, $\alpha\gamma'-\alpha'\gamma$,\ldots
taken between the corresponding pairs of symbols.

The expressions for the invariants and covariants of the binary quantic may be thus exhibited in a compendious form, which renders apparent the symmetry of the several expressions, and shows the connexion of each covariant with the theory of the linear transformation of the quantic.

\clearpage
\setcounter{page}{269}

\textbf{51.} The foregoing investigation establishes the following general theorem: \textit{For a binary quantic of the order $n$, the number of asyzygetic covariants of the degree $m$ is equal to the coefficient of $x^{m}$ in the development of}
\[
\frac{1}{(1-x)(1-x^{2})\ldots(1-x^{n})}\,
\]
\textit{when $m$ is odd; and to the coefficient of $x^{m}$ in the development of}
\[
\frac{1}{(1-x^{2})(1-x^{3})\ldots(1-x^{n})}
\]
\textit{when $m$ is even.} And hence the number of invariants of the degree $m$ is equal to the coefficient of $x^{m}$ in the development of $1/(1-x^{2})(1-x^{3})\ldots(1-x^{n})$, or is equal to the number of partitions of $m$ into the parts $2$, $3$,\ldots, $n$.

\textbf{52.} I have spoken of a covariant as being a function of the coefficients and variables which possesses the property of covariance; but it is proper to remark that every function possessing this property is not thereby a covariant in the strict sense of the term. A function of the coefficients and variables, which is equal to the same function of the transformed coefficients and variables multiplied by a power of the modulus, is a covariant; but there may be functions which possess a more general property, being equal to the same function of the transformed coefficients and variables multiplied by a power of the modulus and also by a factor depending on the coefficients of substitution. Such functions might be termed ``semicovariants''; they correspond to the seminvariants above mentioned.

\textbf{53.} The general expression for the number of the asyzygetic covariants of the degree $m$ of a binary quantic of the order $n$ may be written under the form
\[
\frac{1}{2\pi i}\int \frac{dz}{z^{m+1}(1-z)(1-xz)\ldots(1-x^{n}z)}
\]
evaluated for $x=1$, when $m$ is odd; and under the form
\[
\frac{1}{2\pi i}\int \frac{dz}{z^{m+1}(1-z)(1-x^{2}z)\ldots(1-x^{n}z)}
\]
evaluated for $x=1$, when $m$ is even; or if we please, the two cases may be united in the single formula
\[
\frac{1}{2\pi i}\int \frac{dz}{z^{m+1}(1-xz)(1-x^{2}z)\ldots(1-x^{n}z)}
\]
evaluated for $x=1$.

\clearpage
\setcounter{page}{270}

\textbf{54.} The present Memoir contains the development and establishment of the general theory of the covariants of a binary quantic, so far as relates to the enumeration and determination of the covariants. I have shown that the number of asyzygetic covariants of a given degree can be expressed by means of a definite integral, and that the actual determination of the covinats can be effected by a process which depends upon the solution of a system of partial differential equations. I have also established the connexion between the theory of covariants and the theory of the partition of numbers, and have shown how the general formul\ae\ may be applied to the particular cases of the cubic, quartic, and quintic.

\textbf{55.} It remains to consider the theory of the covariants of quantics of higher orders. The sextic, septimic, octavic, and higher quantics present each a separate problem, the complete solution of which would require a detailed investigation of the same nature as that which has been given for the cubic, quartic, and quintic. The difficulties of such investigations increase rapidly with the order of the quantic; and it is to be observed that the methods which have been employed in the present Memoir are for the most part applicable only to the case of binary quantics. The theory of ternary quantics, and of quantics in general, requires the development of new methods, and the introduction of new principles, which I reserve for future investigation.

\textbf{56.} I will conclude with a remark on the general theory. The covariants of a quantic may be divided into two classes, which may be termed ``absolute covariants'' and ``relative covariants.'' An absolute covariant is a covariant which is equal to the same function of the transformed coefficients and variables; a relative covariant is one which is equal to the same function of the transformed coefficients and variables multiplied by a power of the modulus. In the case of binary quantics, every covariant is a relative covariant; but in the case of ternary and higher quantics, there exist covariants which are absolute covariants. The theory of these absolute covariants presents many interesting questions, which I hope to discuss in a future communication.

\clearpage
\setcounter{page}{271}

\textbf{57.} The theory of the invariants of a binary quantic is intimately connected with the theory of the partition of numbers. In fact, the number of asyzygetic invariants of the degree $m$ of a binary quantic of the order $n$ is equal to the number of partitions of $m$ into the parts $2$, $3$,\ldots, $n$; and the generating function for the asyzygetic invariants is $1/(1-x^{2})(1-x^{3})\ldots(1-x^{n})$. This connexion may be exhibited in a more general form; for if we consider the generating function
\[
\frac{1}{(1-x)(1-xz)(1-xz^{2})\ldots(1-xz^{n})},
\]
then the coefficient of $x^{m}z^{p}$ in the development of this function is equal to the number of asyzygetic covariants of the degree $m$ and the order $p$.

\textbf{58.} The results obtained in the present Memoir may be summed up as follows. For a binary quantic of the order $n$:
\begin{enumerate}[label=(\arabic*)]
\item The number of asyzygetic covariants of the degree $m$ is equal to the coefficient of $x^{m}$ in the development of
\[
\frac{1}{(1-x)(1-x^{2})\ldots(1-x^{n})}
\]
when $m$ is odd; and to the coefficient of $x^{m}$ in the development of
\[
\frac{1}{(1-x^{2})(1-x^{3})\ldots(1-x^{n})}
\]
when $m$ is even.
\item The number of invariants of the degree $m$ is equal to the number of partitions of $m$ into the parts $2$, $3$,\ldots, $n$.
\item The generating function for the covariants is
\[
\frac{1}{(1-x)(1-xz)(1-xz^{2})\ldots(1-xz^{n})}\,.
\]
\item The actual determination of the covariants can be effected by the solution of the differential equations which express the condition of covariance.
\end{enumerate}

\textbf{59.} The foregoing results have been applied in detail to the cases of the cubic, quartic, and quintic; and the complete systems of irreducible covariants and invariants have been obtained for each of these cases. For the cubic, the irreducible covariants are $U$, $H$, $\Phi$, $\nabla$, connected by the syzygy of the degree $6$; the only irreducible invariant is $\nabla$. For the quartic, the irreducible covariants are $U$, $I$, $H$, $J$, $\Phi$, connected by the syzygy of the degree $6$; the only irreducible invariants are $I$, $J$. For the quintic, there are in all twenty-three irreducible covariants, which have been enumerated in the text; and there are four irreducible invariants, of the degrees $4$, $8$, $12$, $18$.

\clearpage
\setcounter{page}{272}

\textbf{60.} I have also considered the general theory of the semicovariants and seminvariants, and have shown that there exist functions which possess a more general property of covariance than that which belongs to covariants in the strict sense. These functions are equal to the same function of the transformed coefficients and variables multiplied by a power of the modulus and also by a factor depending on the coefficients of the substitution. The theory of these functions presents many interesting questions, which I do not however propose to discuss further in the present Memoir.

\textbf{61.} The annexed Tables Nos.\ 23 to 39 contain the developed expressions of the covariants of the binary quintic, with the literal coefficients (these Tables are in fact those which are given in my ``Second Memoir on Quantics,'' Phil.\ Trans.\ vol.\ CXL VI for the year 1856). I give also a Table of the leading coefficients of the several covariants, which serves to fix their numerical values.

\begin{center}
\textsc{Covariants of the Binary Quintic.}
\end{center}

\begin{center}
\textbf{K.\quad No.\ 23.}
\end{center}

\begin{small}
\begin{center}
\begin{tabular}{|r@{\hspace{0.5em}}c@{\hspace{0.5em}}r|r@{\hspace{0.5em}}c@{\hspace{0.5em}}r|r@{\hspace{0.5em}}c@{\hspace{0.5em}}r|r@{\hspace{0.5em}}c@{\hspace{0.5em}}r|}
\hline
$a^{3}bf^{2}$ & \ldots & & $a^{2}cf^{2}$ & $+$ & $1$ & $a^{2}df^{2}$ & $-$ & $1$ & $a^{2}ef^{2}$ & \ldots &\\
$a^{2}cef$ & $+$ & $1$ & $a^{2}def$ & $-$ & $5$ & $a^{2}e^{2}f$ & $+$ & $1$ & $abdf^{2}$ & $-$ & $1$\\
$a^{2}d^{2}f$ & $-$ & $3$ & $a^{2}e^{2}$ & $+$ & $4$ & $abcf^{2}$ & $+$ & $5$ & $abe^{2}f$ & $+$ & $1$\\
$a^{2}de^{2}$ & $+$ & $2$ & $ab^{2}f^{2}$ & $-$ & $1$ & $abdef$ & $-$ & $8$ & $ac^{2}f^{2}$ & $+$ & $3$\\
$ab^{2}ef$ & $-$ & $1$ & $abcef$ & $+$ & $8$ & $abe^{3}$ & $+$ & $3$ & $acdef$ & $-$ & $14$\\
$abcdf$ & $+$ & $14$ & $abd^{2}f$ & $+$ & $11$ & $ac^{2}ef$ & $-$ & $11$ & $ace^{3}$ & $+$ & $8$\\
$abce^{2}$ & $-$ & $11$ & $abde^{2}$ & $-$ & $17$ & $acd^{2}f$ & $+$ & $11$ & $ad^{3}f$ & $+$ & $9$\\
$abd^{2}e$ & $-$ & $1$ & $ac^{2}df$ & $-$ & $11$ & $acde^{2}$ & $+$ & $6$ & $ad^{2}e^{2}$ & $-$ & $6$\\
$ac^{3}f$ & $-$ & $9$ & $ac^{2}e^{2}$ & $-$ & $16$ & $ad^{3}e$ & $-$ & $6$ & $b^{2}cf^{2}$ & $-$ & $2$\\
$ac^{2}de$ & $+$ & $14$ & $acd^{3}$ & $+$ & $44$ & $b^{2}f^{2}$ & $-$ & $4$ & $b^{2}def$ & $+$ & $11$\\
$acd^{3}$ & $-$ & $6$ & $ad^{4}$ & $-$ & $18$ & $b^{2}cef$ & $+$ & $17$ & $b^{2}e^{3}$ & $-$ & $9$\\
$b^{3}df$ & $-$ & $8$ & $b^{2}cf$ & $-$ & $3$ & $b^{2}d^{2}f$ & $+$ & $16$ & $bc^{2}ef$ & $+$ & $1$\\
$b^{2}e^{2}$ & $+$ & $9$ & $b^{2}cdf$ & $-$ & $6$ & $b^{2}de^{2}$ & $-$ & $21$ & $bcd^{2}f$ & $-$ & $14$\\
$b^{2}cf$ & $+$ & $6$ & $b^{2}ce^{2}$ & $+$ & $21$ & $bc^{2}df$ & $-$ & $44$ & $bcde^{2}$ & $+$ & $16$\\
$b^{2}cde$ & $-$ & $16$ & $b^{2}d^{2}e$ & $-$ & $5$ & $bc^{2}e^{2}$ & $+$ & $5$ & $bd^{3}e$ & $-$ & $3$\\
$b^{2}d^{3}$ & $+$ & $8$ & $bc^{2}f$ & $+$ & $6$ & $bcd^{2}e$ & $+$ & $39$ & $c^{3}df$ & $+$ & $6$\\
$bc^{3}e$ & $+$ & $3$ & $bc^{2}de$ & $-$ & $39$ & $bd^{4}$ & $-$ & $12$ & $c^{3}e^{2}$ & $-$ & $8$\\
$bc^{2}d^{2}$ & $-$ & $2$ & $bcd^{3}$ & $+$ & $22$ & $c^{4}f$ & $+$ & $18$ & $c^{2}d^{2}e$ & $+$ & $2$\\
$c^{4}d$ & \ldots & & $c^{4}e$ & $+$ & $12$ & $c^{3}de$ & $-$ & $22$ & $cd^{4}$ & \ldots &\\
\cline{1-9}
\multicolumn{3}{c|}{$\pm\,57$} & \multicolumn{3}{c|}{$\pm\,129$} & \multicolumn{3}{c|}{$\pm\,129$} & \multicolumn{3}{c|}{$\pm\,57$}\\
\hline
\end{tabular}
\end{center}
\end{small}

\hfill $(a,b,c,d,e,f\,\raisebox{0.5ex}{$\pitchfork$}\,x,y)^{5}$

\clearpage
\setcounter{page}{273}

\begin{center}
\textbf{L.\quad No.\ 24.}
\end{center}

\begin{footnotesize}
\begin{center}
\resizebox{\textwidth}{!}{%
\begin{tabular}{|r@{\hspace{0.3em}}c@{\hspace{0.3em}}r|r@{\hspace{0.3em}}c@{\hspace{0.3em}}r|r@{\hspace{0.3em}}c@{\hspace{0.3em}}r|r@{\hspace{0.3em}}c@{\hspace{0.3em}}r|r@{\hspace{0.3em}}c@{\hspace{0.3em}}r|r@{\hspace{0.3em}}c@{\hspace{0.3em}}r|r@{\hspace{0.3em}}c@{\hspace{0.3em}}r|r@{\hspace{0.3em}}c@{\hspace{0.3em}}r|}
\hline
$a^{2}df$ & \ldots & & $a^{2}f^{3}$ & \ldots & & $a^{2}bf^{2}$ & \ldots & & $a^{2}cf^{2}$ & $-$ & $1$ & $a^{2}df^{2}$ & $+$ & $1$ & $a^{2}ef^{2}$ & \ldots & & $a^{2}f^{3}$ & \ldots & & $abf^{3}$ & \ldots &\\
$a^{2}bdf$ & \ldots & & $a^{2}bef$ & \ldots & & $a^{2}cef$ & $-$ & $3$ & $a^{2}def$ & $+$ & $7$ & $a^{2}e^{2}f$ & $-$ & $1$ & $abdf^{2}$ & $+$ & $3$ & $abef^{2}$ & \ldots & & $acef^{2}$ & \ldots &\\
$a^{2}b^{2}e$ & \ldots & & $a^{2}cdf$ & $+$ & $7$ & $a^{2}d^{2}f$ & $+$ & $12$ & $a^{2}e^{3}$ & $-$ & $6$ & $abcf^{2}$ & $-$ & $7$ & $abe^{2}f$ & $-$ & $3$ & $acdf^{2}$ & $-$ & $7$ & $ad^{2}f^{2}$ & $-$ & $2$\\
$a^{2}c^{2}f$ & $+$ & $2$ & $a^{2}ce^{2}$ & $-$ & $10$ & $a^{2}de^{2}$ & $-$ & $9$ & $ab^{2}f^{2}$ & $+$ & $1$ & $abdef$ & $+$ & $26$ & $ac^{2}f^{2}$ & $-$ & $12$ & $ace^{2}f$ & $+$ & $7$ & $ade^{2}f$ & $+$ & $4$\\
$a^{2}cde$ & $-$ & $5$ & $a^{2}d^{2}e$ & $+$ & $3$ & $ab^{2}ef$ & $+$ & $3$ & $abcef$ & $-$ & $26$ & $abe^{3}$ & $-$ & $19$ & $acdef$ & $+$ & $18$ & $ad^{2}ef$ & $+$ & $7$ & $ae^{4}$ & $-$ & $2$\\
$a^{2}d^{3}$ & $+$ & $3$ & $ab^{2}df$ & $-$ & $7$ & $abcdf$ & $-$ & $18$ & $abde^{2}$ & $-$ & $8$ & $ac^{2}df$ & $+$ & $32$ & $ac^{2}e^{2}$ & $+$ & $6$ & $ad^{3}f$ & $-$ & $7$ & $b^{2}ef^{2}$ & \ldots &\\
$ab^{3}f$ & $-$ & $4$ & $ab^{2}e^{2}$ & $+$ & $10$ & $abce^{2}$ & $-$ & $18$ & $abd^{2}f$ & $+$ & $30$ & $ac^{2}ef$ & $-$ & $53$ & $ad^{2}e^{2}$ & $-$ & $15$ & $b^{2}df^{2}$ & $+$ & $10$ & $bcd^{2}$ & $+$ & $5$\\
$ab^{2}de$ & $+$ & $5$ & $abc^{2}f$ & $-$ & $7$ & $abd^{2}e$ & $+$ & $30$ & $ac^{2}df$ & $-$ & $18$ & $acde^{2}$ & $+$ & $53$ & $ad^{3}e$ & $-$ & $39$ & $b^{2}e^{2}f$ & $-$ & $10$ & $bce^{2}f$ & $-$ & $5$\\
$abce^{2}$ & $+$ & $5$ & $abc^{2}de$ & $-$ & $8$ & $ac^{2}f$ & $-$ & $3$ & $ac^{2}e^{2}$ & $+$ & $52$ & $ad^{3}f$ & $-$ & $39$ & $b^{2}f^{2}$ & $+$ & $9$ & $b^{2}def$ & $+$ & $18$ & $b d^{2}$ & $+$ & $5$\\
$abc^{2}d$ & $-$ & $7$ & $abd^{3}$ & $+$ & $9$ & $ac^{2}de$ & $+$ & $45$ & $acd^{2}e$ & $+$ & $6$ & $ad^{4}$ & $-$ & $6$ & $b^{2}cef$ & $+$ & $8$ & $b^{2}e^{3}$ & $-$ & $27$ & $bc^{2}f$ & $-$ & $2$\\
$acd^{2}$ & $+$ & $1$ & $ac^{2}e$ & $+$ & $22$ & $ac^{2}d^{2}$ & $-$ & $39$ & $b^{2}f^{2}$ & $+$ & $6$ & $b^{2}def$ & $+$ & $18$ & $b^{2}e^{3}$ & $-$ & $27$ & $bce^{3}$ & $-$ & $2$ & $c^{2}f^{2}$ & $-$ & $3$\\
$b^{3}f$ & $+$ & $2$ & $ac^{2}d^{2}$ & $-$ & $19$ & $b^{2}df$ & $-$ & $6$ & $b^{2}ef$ & $+$ & $19$ & $b^{2}d^{2}f$ & $-$ & $6$ & $bc^{2}ef$ & $-$ & $30$ & $bd^{2}f$ & $-$ & $22$ & $c^{2}def$ & $+$ & $7$\\
$b^{2}ce$ & $-$ & $5$ & $b^{2}cf$ & $+$ & $7$ & $b^{2}e^{2}$ & $+$ & $27$ & $b^{2}cdf$ & $-$ & $53$ & $b^{2}de^{2}$ & $-$ & $20$ & $bcd^{2}f$ & $-$ & $45$ & $bd^{2}e^{2}$ & $+$ & $19$ & $c^{2}e^{3}$ & $+$ & $2$\\
$b^{2}d^{2}$ & $-$ & $2$ & $b^{2}de$ & $+$ & $2$ & $b^{2}cf$ & $+$ & $15$ & $b^{2}ce^{2}$ & $+$ & $20$ & $bc^{2}df$ & $+$ & $45$ & $bcde^{2}$ & $+$ & $87$ & $c^{2}ef$ & $-$ & $9$ & $cd^{2}f$ & $-$ & $1$\\
$b^{2}cd$ & $+$ & $8$ & $b^{2}ce$ & $-$ & $19$ & $b^{2}cde$ & $-$ & $87$ & $b^{2}d^{2}e$ & $-$ & $25$ & $bc^{2}e^{2}$ & $+$ & $25$ & $bd^{3}e$ & $-$ & $12$ & $c^{2}d^{2}f$ & $+$ & $19$ & $cd^{2}e^{2}$ & $-$ & $8$\\
$bc^{3}$ & $-$ & $3$ & $b^{2}cd^{2}$ & $-$ & $11$ & $b^{2}d^{3}$ & $+$ & $6$ & $bc^{2}f$ & $+$ & $39$ & $b^{2}d^{2}e$ & $-$ & $52$ & $c^{2}df$ & $+$ & $39$ & $c^{2}de^{2}$ & $+$ & $11$ & $d^{3}e$ & $+$ & $3$\\
\multicolumn{1}{|c}{} & & & $bc^{3}d$ & $+$ & $33$ & $bc^{2}e$ & $+$ & $12$ & $bcd^{3}$ & $-$ & $45$ & $bd^{4}$ & \ldots & & $c^{2}e^{2}$ & $-$ & $6$ & $cd^{3}e$ & $-$ & $33$ & \multicolumn{1}{c|}{} & &\\
\multicolumn{1}{|c}{} & & & $c^{4}$ & $-$ & $12$ & $c^{2}d^{2}$ & $+$ & $57$ & $bcd^{2}$ & $+$ & $65$ & $c^{4}f$ & $+$ & $39$ & $c^{2}d^{2}e$ & $-$ & $57$ & $d^{3}$ & $+$ & $12$ & \multicolumn{1}{c|}{} & &\\
\multicolumn{1}{|c}{} & & & \multicolumn{1}{c}{} & & & $c^{4}d$ & $-$ & $24$ & $c^{2}de$ & \ldots & & $c^{2}d^{3}e$ & $-$ & $65$ & $cd^{4}$ & $+$ & $24$ & \multicolumn{1}{c}{} & & & \multicolumn{1}{c|}{} & &\\
\hline
\multicolumn{3}{|c|}{$\pm\,26$} & \multicolumn{3}{c|}{$\pm\,93$} & \multicolumn{3}{c|}{$\pm\,207$} & \multicolumn{3}{c|}{$\pm\,241$} & \multicolumn{3}{c|}{$\pm\,241$} & \multicolumn{3}{c|}{$\pm\,207$} & \multicolumn{3}{c|}{$\pm\,93$} & \multicolumn{3}{c|}{$\pm\,26$}\\
\hline
\end{tabular}%
}
\end{center}
\end{footnotesize}

\hfill $(a,b,c,d,e,f\,\raisebox{0.5ex}{$\pitchfork$}\,x,y)^{7}$

\vspace{1ex}
\begin{center}
\begin{tabular}{r@{\ }l}
No.\ 25, & $Q = +\,1\,a^{2}cdf^{3} + \&c.$\\
No.\ 26, & $Q' = +\,1\,a^{4}f^{4}\phantom{f} + \&c.$
\end{tabular}
\hspace{2em}\raisebox{0.5ex}{$\big\rbrace$} [see \textit{post}, 143.]
\end{center}

\clearpage
\setcounter{page}{274}

\begin{center}
\textsc{G.\quad No.\ 27.}
\end{center}

\begin{small}
\begin{center}
\begin{tabular}{|r@{\hspace{0.4em}}c@{\hspace{0.4em}}r|r@{\hspace{0.4em}}c@{\hspace{0.4em}}r|r@{\hspace{0.4em}}c@{\hspace{0.4em}}r|r@{\hspace{0.4em}}c@{\hspace{0.4em}}r|r@{\hspace{0.4em}}c@{\hspace{0.4em}}r|}
\hline
$a^{3}f^{2}$ & \ldots & & $a^{2}bf^{2}$ & \ldots & & $a^{2}cf^{2}$ & \ldots & & $a^{2}df^{2}$ & \ldots & & $a^{2}ef^{2}$ & \ldots &\\
$a^{3}cef$ & \ldots & & $a^{3}d^{2}f$ & \ldots & & $a^{2}cef$ & \ldots & & $a^{2}def$ & \ldots & & $a^{2}e^{2}f$ & \ldots &\\
$a^{2}c^{2}f$ & \ldots & & $a^{2}cde$ & \ldots & & $a^{2}d^{2}e$ & \ldots & & $ab^{2}f^{2}$ & \ldots & & $abcf^{2}$ & \ldots &\\
$a^{2}cd^{2}$ & \ldots & & $ab^{2}ef$ & \ldots & & $abcdf$ & \ldots & & $abcef$ & \ldots & & $abde^{2}$ & \ldots &\\
$ab^{3}f$ & \ldots & & $ab^{2}ce$ & \ldots & & $ab^{2}d^{2}$ & \ldots & & $abc^{2}e$ & \ldots & & $abcd^{2}$ & \ldots &\\
$abc^{3}$ & \ldots & & $b^{4}f$ & \ldots & & $b^{3}ce$ & \ldots & & $b^{2}c^{2}d$ & \ldots & & $b^{2}cd^{2}$ & \ldots &\\
\hline
\end{tabular}
\end{center}
\end{small}

\hfill $(a,b,c,d,e\,\raisebox{0.5ex}{$\pitchfork$}\,x,y)^{5}$

\vspace{2ex}
\begin{center}
\textsc{H.\quad No.\ 28.}
\end{center}

\begin{small}
\begin{center}
\begin{tabular}{|r@{\hspace{0.4em}}c@{\hspace{0.4em}}r|r@{\hspace{0.4em}}c@{\hspace{0.4em}}r|r@{\hspace{0.4em}}c@{\hspace{0.4em}}r|r@{\hspace{0.4em}}c@{\hspace{0.4em}}r|r@{\hspace{0.4em}}c@{\hspace{0.4em}}r|}
\hline
$a^{4}f$ & \ldots & & $a^{3}bf^{2}$ & \ldots & & $a^{3}cf^{2}$ & \ldots & & $a^{2}df^{2}$ & \ldots & & $a^{2}ef^{2}$ & \ldots &\\
$a^{3}cef$ & \ldots & & $a^{2}d^{2}f$ & \ldots & & $a^{2}ce^{2}$ & \ldots & & $a^{2}def$ & \ldots & & $a^{2}e^{3}$ & \ldots &\\
$a^{2}c^{2}f$ & \ldots & & $a^{2}cd^{2}$ & \ldots & & $ab^{2}f^{2}$ & \ldots & & $abcef$ & \ldots & & $abe^{3}$ & \ldots &\\
$a^{2}d^{3}$ & \ldots & & $abc^{2}f$ & \ldots & & $abcd^{2}$ & \ldots & & $ac^{3}e$ & \ldots & & $ac^{2}d^{2}$ & \ldots &\\
$ab^{3}f$ & \ldots & & $ab^{2}cd$ & \ldots & & $ac^{3}d$ & \ldots & & $b^{4}f$ & \ldots & & $b^{3}cd$ & \ldots &\\
\hline
\end{tabular}
\end{center}
\end{small}

\hfill $(a,b,c,d,e\,\raisebox{0.5ex}{$\pitchfork$}\,x,y)^{5}$

\vspace{2ex}
\begin{center}
\textsc{I.\quad No.\ 29.}
\end{center}

\begin{small}
\begin{center}
\begin{tabular}{|r@{\hspace{0.4em}}c@{\hspace{0.4em}}r|r@{\hspace{0.4em}}c@{\hspace{0.4em}}r|r@{\hspace{0.4em}}c@{\hspace{0.4em}}r|r@{\hspace{0.4em}}c@{\hspace{0.4em}}r|r@{\hspace{0.4em}}c@{\hspace{0.4em}}r|}
\hline
$a^{2}f^{2}$ & \ldots & & $ab^{2}f^{2}$ & \ldots & & $abcf^{2}$ & \ldots & & $abdf^{2}$ & \ldots &\\
$a^{2}cef$ & \ldots & & $a^{2}def$ & \ldots & & $ab^{2}ef$ & \ldots & & $abcef$ & \ldots &\\
$a^{2}d^{2}f$ & \ldots & & $a^{2}de^{2}$ & \ldots & & $ab^{2}df$ & \ldots & & $abcdf$ & \ldots &\\
$a^{2}c^{2}f$ & \ldots & & $ab^{3}f$ & \ldots & & $ab^{2}ce$ & \ldots & & $abc^{2}e$ & \ldots &\\
$ab^{2}d^{2}$ & \ldots & & $abc^{2}d$ & \ldots & & $b^{4}f$ & \ldots & & $b^{2}c^{2}d$ & \ldots &\\
\hline
\end{tabular}
\end{center}
\end{small}

\hfill $(a,b,c,d,e,f\,\raisebox{0.5ex}{$\pitchfork$}\,x,y)^{6}$

\clearpage
\setcounter{page}{275}

\begin{center}
\textsc{J.\quad No.\ 30.}
\end{center}

\begin{small}
\begin{center}
\begin{tabular}{|r@{\hspace{0.4em}}c@{\hspace{0.4em}}r|r@{\hspace{0.4em}}c@{\hspace{0.4em}}r|r@{\hspace{0.4em}}c@{\hspace{0.4em}}r|r@{\hspace{0.4em}}c@{\hspace{0.4em}}r|}
\hline
$a^{2}f^{2}$ & \ldots & & $ab^{2}f^{2}$ & \ldots & & $abcf^{2}$ & \ldots & & $abdf^{2}$ & \ldots &\\
$a^{2}cef$ & \ldots & & $a^{2}d^{2}f$ & \ldots & & $a^{2}e^{2}f$ & \ldots & & $ab^{2}ef$ & \ldots &\\
$a^{2}def$ & \ldots & & $ab^{2}df$ & \ldots & & $abcdf$ & \ldots & & $abcef$ & \ldots &\\
$abe^{3}$ & \ldots & & $ac^{2}f^{2}$ & \ldots & & $acdef$ & \ldots & & $ad^{3}f$ & \ldots &\\
$b^{2}df$ & \ldots & & $b^{2}e^{2}$ & \ldots & & $bc^{2}f$ & \ldots & & $bcd^{2}$ & \ldots &\\
\hline
\end{tabular}
\end{center}
\end{small}

\hfill $(a,b,c,d,e,f\,\raisebox{0.5ex}{$\pitchfork$}\,x,y)^{6}$

\vspace{2ex}
\begin{center}
\textsc{K.\quad No.\ 31.}
\end{center}

\begin{small}
\begin{center}
\begin{tabular}{|r@{\hspace{0.4em}}c@{\hspace{0.4em}}r|r@{\hspace{0.4em}}c@{\hspace{0.4em}}r|r@{\hspace{0.4em}}c@{\hspace{0.4em}}r|r@{\hspace{0.4em}}c@{\hspace{0.4em}}r|}
\hline
$a^{2}f^{2}$ & \ldots & & $ab^{2}f^{2}$ & \ldots & & $abcf^{2}$ & \ldots & & $abdf^{2}$ & \ldots &\\
$a^{2}cef$ & \ldots & & $a^{2}def$ & \ldots & & $ab^{2}ef$ & \ldots & & $abcef$ & \ldots &\\
$a^{2}d^{2}f$ & \ldots & & $ab^{2}df$ & \ldots & & $abcdf$ & \ldots & & $ac^{2}f^{2}$ & \ldots &\\
$acdef$ & \ldots & & $ad^{3}f$ & \ldots & & $b^{2}cf^{2}$ & \ldots & & $b^{2}def$ & \ldots &\\
$bc^{2}ef$ & \ldots & & $bcd^{2}f$ & \ldots & & $c^{3}df$ & \ldots & & $c^{3}e^{2}$ & \ldots &\\
\hline
\end{tabular}
\end{center}
\end{small}

\hfill $(a,b,c,d,e,f\,\raisebox{0.5ex}{$\pitchfork$}\,x,y)^{6}$

\vspace{2ex}
\begin{center}
\textsc{L.\quad No.\ 32.}
\end{center}

\begin{small}
\begin{center}
\begin{tabular}{|r@{\hspace{0.4em}}c@{\hspace{0.4em}}r|r@{\hspace{0.4em}}c@{\hspace{0.4em}}r|r@{\hspace{0.4em}}c@{\hspace{0.4em}}r|r@{\hspace{0.4em}}c@{\hspace{0.4em}}r|r@{\hspace{0.4em}}c@{\hspace{0.4em}}r|}
\hline
$a^{2}f^{2}$ & \ldots & & $ab^{2}f^{2}$ & \ldots & & $abcf^{2}$ & \ldots & & $abdf^{2}$ & \ldots & & $abef^{2}$ & \ldots &\\
$a^{2}cef$ & \ldots & & $a^{2}d^{2}f$ & \ldots & & $a^{2}e^{2}f$ & \ldots & & $ab^{2}ef$ & \ldots &\\
$a^{2}def$ & \ldots & & $ab^{2}df$ & \ldots & & $abcdf$ & \ldots & & $abcef$ & \ldots &\\
$abe^{3}$ & \ldots & & $ac^{2}f^{2}$ & \ldots & & $acdef$ & \ldots & & $ad^{3}f$ & \ldots &\\
$b^{2}df$ & \ldots & & $b^{2}e^{2}$ & \ldots & & $bc^{2}f$ & \ldots & & $bcd^{2}$ & \ldots &\\
\hline
\end{tabular}
\end{center}
\end{small}

\hfill $(a,b,c,d,e,f\,\raisebox{0.5ex}{$\pitchfork$}\,x,y)^{6}$

\clearpage
\setcounter{page}{276}

\begin{center}
\textsc{M.\quad No.\ 33.}
\end{center}

\begin{small}
\begin{center}
\begin{tabular}{|r@{\hspace{0.4em}}c@{\hspace{0.4em}}r|r@{\hspace{0.4em}}c@{\hspace{0.4em}}r|r@{\hspace{0.4em}}c@{\hspace{0.4em}}r|r@{\hspace{0.4em}}c@{\hspace{0.4em}}r|}
\hline
$a^{3}bf$ & \ldots & & $a^{3}cf$ & \ldots & & $a^{3}df$ & \ldots & & $a^{3}ef$ & \ldots &\\
$a^{3}ce$ & \ldots & & $a^{3}d^{2}$ & \ldots & & $a^{2}b^{2}f$ & \ldots & & $a^{2}bcf$ & \ldots &\\
$a^{2}bdf$ & \ldots & & $a^{2}be^{2}$ & \ldots & & $a^{2}c^{2}f$ & \ldots & & $a^{2}cde$ & \ldots &\\
$a^{2}d^{3}$ & \ldots & & $ab^{3}f$ & \ldots & & $ab^{2}ce$ & \ldots & & $ab^{2}d^{2}$ & \ldots &\\
$abc^{2}d$ & \ldots & & $ab^{2}c^{2}$ & \ldots & & $b^{4}f$ & \ldots & & $b^{3}ce$ & \ldots &\\
$b^{2}c^{2}d$ & \ldots & & $b^{2}cd^{2}$ & \ldots & & $bc^{4}$ & \ldots &\\
\hline
\end{tabular}
\end{center}
\end{small}

\hfill $(a,b,c,d,e\,\raisebox{0.5ex}{$\pitchfork$}\,x,y)^{4}$

\vspace{2ex}
\begin{center}
\textsc{N.\quad No.\ 34.}
\end{center}

\begin{small}
\begin{center}
\begin{tabular}{|r@{\hspace{0.4em}}c@{\hspace{0.4em}}r|r@{\hspace{0.4em}}c@{\hspace{0.4em}}r|r@{\hspace{0.4em}}c@{\hspace{0.4em}}r|r@{\hspace{0.4em}}c@{\hspace{0.4em}}r|}
\hline
$a^{3}bf$ & \ldots & & $a^{3}cf$ & \ldots & & $a^{3}df$ & \ldots & & $a^{3}ef$ & \ldots &\\
$a^{2}b^{2}f$ & \ldots & & $a^{2}bcf$ & \ldots & & $a^{2}bdf$ & \ldots & & $a^{2}be^{2}$ & \ldots &\\
$a^{2}c^{2}f$ & \ldots & & $a^{2}cde$ & \ldots & & $a^{2}d^{3}$ & \ldots & & $ab^{3}f$ & \ldots &\\
$ab^{2}ce$ & \ldots & & $ab^{2}d^{2}$ & \ldots & & $abc^{2}d$ & \ldots & & $ab^{2}c^{2}$ & \ldots &\\
$b^{4}f$ & \ldots & & $b^{3}ce$ & \ldots & & $b^{2}c^{2}d$ & \ldots & & $b^{2}cd^{2}$ & \ldots &\\
\hline
\end{tabular}
\end{center}
\end{small}

\hfill $(a,b,c,d,e\,\raisebox{0.5ex}{$\pitchfork$}\,x,y)^{4}$

\vspace{2ex}
\begin{center}
\textsc{O.\quad No.\ 35.}
\end{center}

\begin{small}
\begin{center}
\begin{tabular}{|r@{\hspace{0.4em}}c@{\hspace{0.4em}}r|r@{\hspace{0.4em}}c@{\hspace{0.4em}}r|r@{\hspace{0.4em}}c@{\hspace{0.4em}}r|r@{\hspace{0.4em}}c@{\hspace{0.4em}}r|}
\hline
$a^{2}bf$ & \ldots & & $a^{2}cf$ & \ldots & & $a^{2}df$ & \ldots & & $a^{2}ef$ & \ldots &\\
$ab^{2}f$ & \ldots & & $abcf$ & \ldots & & $abdf$ & \ldots & & $abe^{2}$ & \ldots &\\
$ac^{2}f$ & \ldots & & $acde$ & \ldots & & $ad^{3}$ & \ldots & & $b^{3}f$ & \ldots &\\
$b^{2}ce$ & \ldots & & $b^{2}d^{2}$ & \ldots & & $bc^{2}d$ & \ldots & & $c^{4}$ & \ldots &\\
\hline
\end{tabular}
\end{center}
\end{small}

\hfill $(a,b,c,d,e,f\,\raisebox{0.5ex}{$\pitchfork$}\,x,y)^{5}$

\vspace{2ex}
\begin{center}
\textsc{P.\quad No.\ 36.}
\end{center}

\begin{small}
\begin{center}
\begin{tabular}{|r@{\hspace{0.4em}}c@{\hspace{0.4em}}r|r@{\hspace{0.4em}}c@{\hspace{0.4em}}r|r@{\hspace{0.4em}}c@{\hspace{0.4em}}r|r@{\hspace{0.4em}}c@{\hspace{0.4em}}r|}
\hline
$a^{2}bf$ & \ldots & & $a^{2}cf$ & \ldots & & $a^{2}df$ & \ldots &\\
$ab^{2}f$ & \ldots & & $abcf$ & \ldots & & $abdf$ & \ldots &\\
$ac^{2}f$ & \ldots & & $acde$ & \ldots & & $ad^{3}$ & \ldots &\\
$b^{3}f$ & \ldots & & $b^{2}ce$ & \ldots & & $b^{2}d^{2}$ & \ldots &\\
$bc^{2}d$ & \ldots & & $c^{4}$ & \ldots &\\
\hline
\end{tabular}
\end{center}
\end{small}

\hfill $(a,b,c,d,e,f\,\raisebox{0.5ex}{$\pitchfork$}\,x,y)^{5}$

\vspace{2ex}
\begin{center}
\begin{tabular}{r@{\ }l}
No.\ 37, & $A = +\,1\,a^{3}cde + \&c.$\\
No.\ 38, & $B = +\,1\,a^{2}b^{2}de + \&c.$\\
No.\ 39, & $C = +\,1\,a^{2}bc^{2}e + \&c.$
\end{tabular}
\end{center}

\vspace{1ex}
\hfill $35\text{---}2$



% =====================================================================
% PAPER [142] --- NUMERICAL TABLES SUPPLEMENTARY TO SECOND MEMOIR ON QUANTICS
% =====================================================================

\clearpage
\setcounter{page}{277}

\begin{center}
\textsc{Numerical Tables Supplementary to Second Memoir on Quantics.}
\end{center}

\begin{center}
[Now first published (1889).]
\end{center}

In the present paper I arrange in a more compendious form and continue to a much greater extent the tables (first of each pair) given Nos.\ 35---39 of my Second Memoir on Quantics, 141, pp.\ 260---264, which relate to the cubic, the quartic and the quintic functions; and I give the like tables for the sextic, the septimic and the octavic functions respectively. The cubic table exhibits the coefficients of the several $xz$ terms of the function $1 \div (1-z.1-xz.1-x^{2}z.1-x^{3}z)$, or, what is the same thing, it gives the number of partitions of a given number into a given number of parts, the parts being $0$, $1$, $2$, $3$, (repetitions admissible): or again, regarding the letters $a$, $b$, $c$, $d$, as having the weights $0$, $1$, $2$, $3$ respectively, it shows the number of literal terms of a given degree and given weight. And similarly for the quartic, quintic, sextic, septimic and octavic tables respectively, the parts of course being $0$, $1$, \ldots\ up to $4$, $5$, $6$, $7$ or $8$, and the letters being $a$, $b$, \ldots\ up to $e$, $f$, $g$, $h$ or $i$. The extent of the tables is as follows:

\begin{center}
\begin{tabular}{l@{\hspace{2em}}l@{\hspace{2em}}l}
cubic table extends to deg-weight & $18$---$27$ & \\
quartic   & \phantom{ext} & $18$---$36$ \\
quintic   & \phantom{ext} & $18$---$45$ \\
sextic    & \phantom{ext} & $15$---$45$ \\
septimic  & \phantom{ext} & $12$---$42$ \\
octavic   & \phantom{ext} & $10$---$40$ \\
\end{tabular}
\end{center}

viz.\ for the quintic, the sextic and the octavic functions these are the deg-weights of the highest invariants respectively. I designate the Tables as the $ad$-, $ae$-, $af$-, $ag$-, $ah$- and $ai$-tables respectively.

It is to be noticed that in the several tables the lower part of each column is for shortness omitted; the column has to be completed by taking into it the series of bottom terms of each of the preceding columns: thus in the $af$- or quintic table the complete column for degree $3$ would be

\begin{center}
\begin{tabular}{r|r}
D & 3 \\
W & $8$-$7$ \\ \hline
$-0$ & $6$ \\
$-1$ & $6$ \\
$-2$ & $5$ \\
$-3$ & $4$ \\
$-4$ & $3$ \\
$-5$ & $2$ \\
$-6$ & $1$ \\
$-7$ & $1$ \\
\end{tabular}
\end{center}

where the concluding terms $2$, $1$, $1$ are the bottom terms of the three preceding columns respectively. And the meaning is that for degree $3$, and weight $8$, or $7$, the number of terms is $=6$; for weight $7-1$, $=6$, the number of terms is $=6$; and similarly for weights $5$, $4$, $3$, $2$, $1$, $0$ the numbers are $5$, $4$, $3$, $2$, $1$, $1$; the numbers are those of the terms

\begin{center}
\begin{tabular}{c|ccccccccc}
W. & $0$ & $1$ & $2$ & $3$ & $4$ & $5$ & $6$ & $7$ & $8$ \\ \hline
& $a^{3}$ & $a^{2}b$ & $a^{2}c$ & $a^{2}d$ & $a^{2}e$ & $a^{2}f$ & $abf$ & $acf$ & $adf$ \\
& & $ab^{2}$ & $abc$ & $abd$ & $abe$ & $ace$ & $ade$ & $ae^{2}$ & \\
& & & $b^{3}$ & $ac^{2}$ & $acd$ & $ad^{2}$ & $b^{2}f$ & $bef$ & \\
& & & & $b^{2}c$ & $b^{2}d$ & $b^{2}e$ & $bce$ & $bde$ & \\
& & & & & $bc^{2}$ & $bcd$ & $bd^{2}$ & $c^{2}e$ & \\
& & & & & & $c^{3}$ & $c^{2}d$ & $cd^{2}$ & \\
\hline
No. & $1$ & $1$ & $2$ & $3$ & $4$ & $5$ & $6$ & $6$ & $6$ \\
\end{tabular}
\end{center}

The like remarks and explanations apply to the other tables.

\begin{center}
\textsc{$ad$-Table.}
\end{center}

\begin{center}
\begin{small}
\resizebox{\textwidth}{!}{%
\begin{tabular}{r|*{19}{r}}
D & $0$ & $1$ & $2$ & $3$ & $4$ & $5$ & $6$ & $7$ & $8$ & $9$ & $10$ & $11$ & $12$ & $13$ & $14$ & $15$ & $16$ & $17$ & $18$ \\ \hline
W & $0$ & $2$-$1$ & $3$ & $5$-$4$ & $6$ & $8$-$7$ & $9$ & $11$-$10$ & $12$ & $14$-$13$ & $15$ & $17$-$16$ & $18$ & $20$-$19$ & $21$ & $23$-$22$ & $24$ & $26$-$25$ & $27$ \\[4pt] \hline
$-0$ & $1$ & $1$ & $2$ & $3$ & $5$ & $6$ & $8$ & $10$ & $13$ & $15$ & $18$ & $21$ & $25$ & $28$ & $32$ & $36$ & $41$ & $45$ & $50$ \\
$-1$ & & & $2$ & $3$ & $4$ & $6$ & $8$ & $10$ & $12$ & $15$ & $18$ & $21$ & $24$ & $28$ & $32$ & $36$ & $40$ & $45$ & $50$ \\
$-2$ & & & & & $4$ & $5$ & $7$ & $9$ & $12$ & $14$ & $17$ & $20$ & $24$ & $27$ & $31$ & $35$ & $40$ & $44$ & $49$ \\
$-3$ & & & & & & & $7$ & $8$ & $11$ & $13$ & $17$ & $19$ & $23$ & $26$ & $31$ & $34$ & $39$ & $43$ & $49$ \\
$-4$ & & & & & & & & & $10$ & $12$ & $15$ & $18$ & $22$ & $25$ & $29$ & $33$ & $38$ & $42$ & $47$ \\
$-5$ & & & & & & & & & & & $14$ & $16$ & $20$ & $23$ & $28$ & $31$ & $36$ & $40$ & $46$ \\
$-6$ & & & & & & & & & & & & & $19$ & $21$ & $26$ & $29$ & $35$ & $38$ & $44$ \\
$-7$ & & & & & & & & & & & & & & & $24$ & $27$ & $32$ & $36$ & $42$ \\
$-8$ & & & & & & & & & & & & & & & & & $30$ & $33$ & $39$ \\
$-9$ & & & & & & & & & & & & & & & & & & & $37$ \\
\end{tabular}%
}
\end{small}
\end{center}

\clearpage
\setcounter{page}{278}

\begin{center}
\textsc{$ae$-Table.}
\end{center}

\begin{center}
\begin{small}
\resizebox{\textwidth}{!}{%
\begin{tabular}{r|*{19}{r}}
D & $1$ & $2$ & $3$ & $4$ & $5$ & $6$ & $7$ & $8$ & $9$ & $10$ & $11$ & $12$ & $13$ & $14$ & $15$ & $16$ & $17$ & $18$ \\ \hline
W & $3$ & $4$ & $5$ & $8$ & $10$ & $13$ & $14$ & $16$ & $18$ & $20$ & $22$ & $24$ & $26$ & $28$ & $30$ & $32$ & $34$ & $36$ \\[4pt] \hline
$-0$ & $1$ & $1$ & $3$ & $5$ & $8$ & $12$ & $18$ & $24$ & $33$ & $43$ & $55$ & $69$ & $86$ & $104$ & $126$ & $150$ & $177$ & $207$ \\
$-1$ & $1$ & $2$ & $4$ & $7$ & $11$ & $16$ & $23$ & $31$ & $41$ & $53$ & $67$ & $83$ & $102$ & $123$ & $147$ & $174$ & $204$ & $237$ \\
$-2$ & $1$ & $2$ & $4$ & $7$ & $11$ & $16$ & $23$ & $31$ & $41$ & $53$ & $67$ & $83$ & $102$ & $123$ & $147$ & $174$ & $204$ & $237$ \\
$-3$ & & $1$ & $3$ & $5$ & $9$ & $14$ & $20$ & $28$ & $38$ & $49$ & $63$ & $79$ & $97$ & $118$ & $142$ & $168$ & $198$ & $231$ \\
$-4$ & & & & $5$ & $8$ & $13$ & $19$ & $27$ & $36$ & $48$ & $61$ & $77$ & $95$ & $116$ & $139$ & $166$ & $195$ & $228$ \\
$-5$ & & & & & & $6$ & $10$ & $16$ & $23$ & $32$ & $43$ & $56$ & $71$ & $89$ & $109$ & $132$ & $158$ & $187$ \\
$-6$ & & & & & & & & $9$ & $14$ & $21$ & $30$ & $40$ & $53$ & $68$ & $85$ & $105$ & $128$ & $153$ \\
$-7$ & & & & & & & & & & $11$ & $17$ & $25$ & $35$ & $47$ & $61$ & $78$ & $97$ & $119$ \\
$-8$ & & & & & & & & & & & & $15$ & $22$ & $32$ & $43$ & $57$ & $73$ & $92$ \\
$-9$ & & & & & & & & & & & & & & $18$ & $26$ & $37$ & $50$ & $65$ \\
\end{tabular}%
}
\end{small}
\end{center}

\clearpage
\setcounter{page}{279}

\begin{center}
\textsc{$af$-Table.}
\end{center}

\begin{center}
\begin{small}
\resizebox{\textwidth}{!}{%
\begin{tabular}{r|*{16}{r}}
D & $1$ & $2$ & $3$ & $4$ & $5$ & $6$ & $7$ & $8$ & $9$ & $10$ & $11$ & $12$ & $13$ & $14$ & $15$ & $16$ \\ \hline
W & $3$ & $6$ & $9$ & $12$ & $15$ & $18$ & $21$ & $24$ & $27$ & $30$ & $33$ & $36$ & $39$ & $42$ & $45$ & $48$ \\[4pt] \hline
$-0$ & $1$ & $1$ & $4$ & $8$ & $18$ & $32$ & $58$ & $94$ & $151$ & $227$ & $338$ & $480$ & $676$ & $920$ \\
$-1$ & $1$ & $3$ & $8$ & $16$ & $32$ & $55$ & $94$ & $146$ & $221$ & $332$ & $480$ & $668$ & $920$ \\
$-2$ & $1$ & $3$ & $7$ & $16$ & $30$ & $55$ & $90$ & $139$ & $217$ & $319$ & $464$ & $648$ & $896$ \\
$-3$ & & $1$ & $2$ & $7$ & $14$ & $29$ & $51$ & $88$ & $134$ & $205$ & $310$ & $446$ & $634$ & $870$ \\
$-4$ & & & & $4$ & $10$ & $23$ & $42$ & $76$ & $123$ & $196$ & $293$ & $431$ & $608$ & $847$ \\
$-5$ & & & & & $3$ & $9$ & $19$ & $39$ & $68$ & $116$ & $182$ & $280$ & $408$ & $587$ & $813$ \\
$-6$ & & & & & & $6$ & $16$ & $32$ & $61$ & $103$ & $169$ & $258$ & $387$ & $553$ & $780$ \\
$-7$ & & & & & & & $5$ & $12$ & $28$ & $52$ & $90$ & $146$ & $241$ & $359$ & $525$ & $737$ \\
$-8$ & & & & & & & & $10$ & $22$ & $46$ & $87$ & $158$ & $272$ & $441$ & $686$ & $1038$ \\
$-9$ & & & & & & & & & $6$ & $17$ & $37$ & $75$ & $136$ & $247$ & $403$ & $663$ \\
\end{tabular}%
}
\end{small}
\end{center}

\clearpage
\setcounter{page}{280}

\begin{center}
\textsc{$ag$-Table.}
\end{center}

\begin{center}
\begin{small}
\resizebox{\textwidth}{!}{%
\begin{tabular}{r|*{16}{r}}
D & $1$ & $2$ & $3$ & $4$ & $5$ & $6$ & $7$ & $8$ & $9$ & $10$ & $11$ & $12$ & $13$ & $14$ & $15$ & $16$ \\ \hline
W & $4$-$5$ & $7$ & $10$-$11$ & $13$-$14$ & $16$-$17$ & $19$-$20$ & $22$-$23$ & $25$-$26$ & $28$-$29$ & $31$-$32$ & $34$-$35$ & $38$-$39$ & $41$-$42$ & $45$ \\[4pt] \hline
$-0$ & $1$ & $1$ & $4$ & $10$ & $24$ & $49$ & $94$ & $169$ & $289$ & $468$ & $734$ & $1117$ & $1656$ \\
$-1$ & & $1$ & $4$ & $10$ & $23$ & $48$ & $94$ & $166$ & $285$ & $464$ & $734$ & $1109$ & $1646$ \\
$-2$ & & $1$ & $3$ & $9$ & $23$ & $46$ & $90$ & $162$ & $282$ & $454$ & $722$ & $1093$ & $1634$ \\
$-3$ & & & $3$ & $8$ & $20$ & $43$ & $88$ & $155$ & $272$ & $441$ & $709$ & $1069$ & $1605$ \\
$-4$ & & & $2$ & $7$ & $19$ & $39$ & $81$ & $146$ & $263$ & $424$ & $686$ & $1038$ & $1572$ \\
$-5$ & & & $2$ & $5$ & $16$ & $35$ & $76$ & $136$ & $247$ & $403$ & $663$ & $1000$ & $1524$ \\
$-6$ & & & & $4$ & $14$ & $30$ & $68$ & $125$ & $233$ & $379$ & $629$ & $957$ & $1475$ \\
$-7$ & & & & $3$ & $11$ & $26$ & $61$ & $112$ & $214$ & $354$ & $598$ & $908$ & $1410$ \\
$-8$ & & & & & $9$ & $21$ & $52$ & $100$ & $197$ & $325$ & $558$ & $856$ & $1346$ \\
$-9$ & & & & & $6$ & $17$ & $46$ & $87$ & $176$ & $297$ & $520$ & $799$ & $1271$ \\
$-10$ & & & & & $5$ & $13$ & $37$ & $75$ & $158$ & $268$ & $477$ & $742$ & $1197$ \\
\end{tabular}%
}
\end{small}
\end{center}

\clearpage
\setcounter{page}{281}

\begin{center}
\textsc{$ah$-Table.}
\end{center}

\begin{center}
\begin{small}
\resizebox{\textwidth}{!}{%
\begin{tabular}{r|*{16}{r}}
D & $1$ & $2$ & $3$ & $4$ & $5$ & $6$ & $7$ & $8$ & $9$ & $10$ & $11$ & $12$ & $13$ & $14$ & $15$ & $16$ \\ \hline
W & $4$ & $8$ & $13$ & $16$ & $20$ & $24$ & $28$ & $32$ & $36$ & $40$ & $44$ & $48$ & $52$ & $56$ & $60$ & $64$ \\[4pt] \hline
$-0$ & $1$ & $1$ & $5$ & $12$ & $33$ & $73$ & $151$ & $289$ & $526$ & $910$ & $1514$ & $1502$ & \ldots \\
$-1$ & $1$ & $4$ & $12$ & $31$ & $70$ & $146$ & $282$ & $515$ & $894$ & $1492$ & \ldots \\
$-2$ & $1$ & $3$ & $11$ & $28$ & $66$ & $139$ & $272$ & $499$ & $873$ & $1460$ & \ldots \\
$-3$ & & $3$ & $10$ & $27$ & $63$ & $134$ & $263$ & $486$ & $851$ & $1430$ & \ldots \\
$-4$ & & $2$ & $8$ & $23$ & $57$ & $123$ & $247$ & $461$ & $816$ & $1379$ & \ldots \\
$-5$ & & $2$ & $7$ & $21$ & $52$ & $116$ & $233$ & $440$ & $783$ & $1331$ & \ldots \\
$-6$ & & & $5$ & $17$ & $45$ & $103$ & $214$ & $409$ & $738$ & $1265$ & \ldots \\
$-7$ & & & $4$ & $15$ & $40$ & $94$ & $197$ & $383$ & $696$ & $1214$ & \ldots \\
$-8$ & & & $3$ & $11$ & $33$ & $81$ & $176$ & $348$ & $645$ & $1127$ & \ldots \\
$-9$ & & & & $9$ & $28$ & $71$ & $158$ & $319$ & $597$ & $1057$ & \ldots \\
\end{tabular}%
}
\end{small}
\end{center}

\clearpage
\setcounter{page}{282}

\begin{center}
\textsc{$ai$-Table.}
\end{center}

\begin{center}
\begin{small}
\resizebox{\textwidth}{!}{%
\begin{tabular}{r|*{16}{r}}
D & $1$ & $2$ & $3$ & $4$ & $5$ & $6$ & $7$ & $8$ & $9$ & $10$ & $11$ & $12$ & $13$ & $14$ & $15$ & $16$ \\ \hline
W & $4$ & $8$ & $13$ & $16$ & $20$ & $24$ & $28$ & $32$ & $36$ & $40$ & $44$ & $48$ & $52$ & $56$ & $60$ & $64$ \\[4pt] \hline
$-0$ & $1$ & $1$ & $5$ & $13$ & $31$ & $73$ & $151$ & $289$ & $526$ & $519$ & $910$ & $902$ & $1514$ & $1502$ & \ldots \\
$-1$ & $1$ & $4$ & $12$ & $31$ & $70$ & $146$ & $282$ & $515$ & $894$ & $1492$ & \ldots \\
$-2$ & $1$ & $3$ & $11$ & $28$ & $66$ & $139$ & $272$ & $499$ & $873$ & $1460$ & \ldots \\
$-3$ & $3$ & $10$ & $27$ & $63$ & $134$ & $263$ & $486$ & $851$ & $1430$ & \ldots \\
$-4$ & $2$ & $8$ & $23$ & $57$ & $123$ & $247$ & $461$ & $816$ & $1379$ & \ldots \\
$-5$ & $2$ & $7$ & $21$ & $52$ & $116$ & $233$ & $440$ & $783$ & $1331$ & \ldots \\
$-6$ & $5$ & $17$ & $45$ & $103$ & $214$ & $409$ & $738$ & $1265$ & \ldots \\
$-7$ & $4$ & $15$ & $40$ & $94$ & $197$ & $383$ & $696$ & $1214$ & \ldots \\
$-8$ & $3$ & $11$ & $33$ & $81$ & $176$ & $348$ & $645$ & $1127$ & \ldots \\
$-9$ & & $9$ & $28$ & $71$ & $158$ & $319$ & $597$ & $1057$ & \ldots \\
\end{tabular}%
}
\end{small}
\end{center}

\vspace{2ex}

The remaining tables continue the numerical values to higher degrees and weights, following the same pattern of triangular arrays. The column headings indicate the degree and weight ranges, and the row labels on the left indicate the weight offset from the maximum.



% =====================================================================
% PAPER [143] --- TABLES OF THE COVARIANTS M TO W OF THE BINARY QUINTIC
% =====================================================================

\clearpage
\setcounter{page}{282}

\begin{center}
\textsc{143.}
\end{center}

\begin{center}
\textsc{Tables of the Covariants M to W of the Binary Quintic: from the Second, Third, Fifth, Eighth, Ninth and Tenth Memoirs on Quantics.}
\end{center}

\begin{center}
[Arranged in the present form, 1889.]
\end{center}

The binary quintic has in all (including the quintic itself and the invariants) 23 covariants, which I have represented by the capital letters, A, B, C, \ldots W (alternative forms of two of these are denoted by Q$'$ and S$'$). The covariants A, \ldots L, and also Q, Q$'$ were given in my Second Memoir on Quantics, and except Q and Q$'$ are reproduced in the present reprint thereof, 141; in all these I gave not only the literal terms actually presenting themselves, but also the terms with zero coefficients; in the other covariants however, or in most of them, the terms with zero coefficients were omitted. It is very desirable to have in every case the complete series of literal terms, and in the covariants as here printed they are accordingly inserted: the number of terms is in each case known beforehand by the foregoing $af$-table, 142, and any omission is thus precluded; by means of this $af$-table we have the numbers of terms as shown in the following list.

I have throughout (as was done in the Ninth and Tenth Memoirs) expressed the literal terms in a slightly different form from that employed in the Second Memoir: this is done in order to show at a glance in each column the set of terms which contain a given power of $a$, and in each such set the terms which contain a given power of $b$.

The numerical verifications are also given not only for the entire column but for each set of terms containing the same power of $a$; viz.\ in most cases, but not always, the positive and negative coefficients of a set have equal sums, which are shown by a number with the sign $\pm$ prefixed. The verification is in some cases given in regard to the subsets involving the same powers of $a$ and $b$, here also the sums of the positive and negative coefficients are not in every case equal. The cases of inequality will be referred to at the end of this paper.

\clearpage
\setcounter{page}{283}

The whole series of covariants is as follows:

\begin{center}
\begin{small}
\begin{tabular}{ccclc}
Mem. & No.\ of table. & & \multicolumn{1}{c}{Covariant} & deg-weight \\
\hline
2 & 13 & A & $=\phantom{(}1,\,1,\,1,\,1,\,1,\,1\between(x,y)^{5}$ & $1$ ($0$ \ldots $5$) \\
\phantom{2} & 14 & B & $=\phantom{(}3,\,3,\,3\between(x,y)^{2}$ & $2$ ($4.\,6$) \\
\phantom{2} & 15 & C & $=(2,\,2,\,3,\,3,\,3,\,2,\,2\between(x,y)^{6}$ & $2$ ($2$ \ldots $8$) \\
\phantom{2} & 16 & D & $=\phantom{(}6,\,6,\,6,\,6\between(x,y)^{3}$ & $3$ ($6$..\,$9$) \\
\phantom{2} & 17 & E & $=\phantom{(}5,\,6,\,6,\,6,\,6,\,5\between(x,y)^{3}$ & $3$ ($5$ \ldots $10$) \\
\phantom{2} & 18 & F & $=(3,\,4,\,5,\,6,\,6,\,6,\,6,\,5,\,4,\,3\between(x,y)^{9}$ & $3$ ($3$ \ldots $12$) \\
\phantom{2} & 19 & G & $=\phantom{(}12\between(x,y)^{0}$, Invt. & $4$-$10$ \\
\phantom{2} & 20 & H & $=(11,\,11,\,12,\,11,\,11\between(x,y)^{4}$ & $4$ ($8$ \ldots $12$) \\
\phantom{2} & 21 & I & $=(9,\,11,\,11,\,12,\,11,\,11,\,9\between(x,y)^{6}$ & $4$ ($7$ \ldots $13$) \\
\phantom{2} & 22 & J & $=\phantom{(}20,\,20\between(x,y)^{1}$ & $5$ ($12$, $13$) \\
\phantom{2} & 23 & K & $=(19,\,20,\,20,\,19\between(x,y)^{3}$ & $5$ ($11$..\,$14$) \\
\phantom{2} & 24 & L & $=(16,\,18,\,19,\,20,\,20,\,19,\,18,\,16\between(x,y)^{7}$ & $5$ ($9$ \ldots $16$) \\
8 & 83 & M & $=\phantom{(}32,\,32,\,32\between(x,y)^{2}$ & $6$ ($14$..\,$16$) \\
\phantom{8} & 84 & N & $=(30,\,32,\,32,\,32,\,30\between(x,y)^{4}$ & $6$ ($13$ \ldots $17$) \\
9 & 90 & O & $=\phantom{(}49,\,49\between(x,y)^{1}$ & $7$ ($17$, $18$) \\
\phantom{9} & 91 & P & $=(46,\,48,\,49,\,49,\,48,\,46\between(x,y)^{5}$ & $7$ ($15$ \ldots $20$) \\
2 & \begin{tabular}{c}Q 25\\ Q$'$ 26\end{tabular} & \begin{tabular}{c}Q\\ Q$'$\end{tabular} & $=\phantom{(}73\between(x,y)^{0}$, Invt. & $8$-$20$ \\
9 & 92 & R & $=\phantom{(}71,\,73,\,71\between(x,y)^{2}$ & $8$ ($19$.\,$21$) \\
9 & \begin{tabular}{c}S 93\\ S 93 bis\end{tabular} & \begin{tabular}{c}S\\ S$'$\end{tabular} & $=(101,\,102,\,102,\,101\between(x,y)^{3}$ & $9$ ($21$..\,$24$) \\
10 & 94 & T & $=\phantom{(}190,\,190\between(x,y)^{1}$ & $11$ ($27$, $28$) \\
3 & 29 & U & $=\phantom{(}252\between(x,y)^{0}$, Invt. & $12$-$30$ \\
9 & 95 & V & $=\phantom{(}325,\,325\between(x,y)^{1}$ & $13$ ($32$, $33$) \\
5 & 29A & W & $=\phantom{(}967\between(x,y)^{0}$, Invt. & $18$-$45$ \\
\hline
\end{tabular}
\end{small}
\end{center}

\hfill $36\text{---}2$

\clearpage
\setcounter{page}{284}

\begin{center}
\textbf{M.\quad No.\ 83.}
\end{center}

\begin{footnotesize}
\begin{center}
\begin{tabular}{|r@{\hspace{0.4em}}c@{\hspace{0.4em}}r|r@{\hspace{0.4em}}c@{\hspace{0.4em}}r|r@{\hspace{0.4em}}c@{\hspace{0.4em}}r|r@{\hspace{0.4em}}c@{\hspace{0.4em}}r|}
\hline
$a^{4}$ $b^{2}df^{2}$ & \ldots & & $a^{3}$ $b^{3}f^{3}$ & \ldots & & $a^{2}$ $b^{4}f^{3}$ & $-$ & $1$ & $a^{2}$ $b^{3}df^{3}$ & $+$ & $1$ \\
$a^{3}$ $b^{2}ef^{2}$ & \ldots & & $a^{2}$ $b^{3}f^{2}$ & \ldots & & $b^{4}df^{2}$ & $+$ & $1$ & $b^{5}ef^{2}$ & $-$ & $1$ \\
$a^{2}$ $c^{2}f^{2}$ & $-$ & $1$ & $a^{2}$ $b^{2}cdf$ & $+$ & $2$ & $c^{2}e^{2}f$ & $+$ & $5$ & $d^{3}f$ & $+$ & $1$ \\
$c^{2}ef$ & $-$ & $2$ & $c^{2}df^{2}$ & $-$ & $2$ & $c^{2}ef$ & $+$ & $2$ & $d^{2}e^{2}$ & $-$ & $1$ \\
$d^{3}f$ & $-$ & $1$ & $cdef$ & $-$ & $6$ & $cd^{2}f$ & $+$ & $6$ & $cde^{2}$ & $-$ & $5$ \\
$d^{2}e^{2}$ & $+$ & $2$ & $ce^{3}$ & $+$ & $3$ & $de^{3}$ & $-$ & $3$ & $e^{4}$ & $+$ & $1$ \\
\hline
\end{tabular}
\end{center}
\end{footnotesize}

\hfill $(a,b,c,d,e,f\between x,y)^{4}$

\vspace{2ex}

\begin{center}
\textbf{N.\quad No.\ 84.}
\end{center}

\begin{footnotesize}
\begin{center}
\begin{tabular}{|r@{\hspace{0.4em}}c@{\hspace{0.4em}}r|r@{\hspace{0.4em}}c@{\hspace{0.4em}}r|r@{\hspace{0.4em}}c@{\hspace{0.4em}}r|r@{\hspace{0.4em}}c@{\hspace{0.4em}}r|r@{\hspace{0.4em}}c@{\hspace{0.4em}}r|}
\hline
$a^{4}$ $bcdf$ & $-$ & $1$ & $a^{3}$ $b^{2}df$ & $-$ & $4$ & $a^{3}$ $b^{2}ef$ & \ldots & & $a^{2}$ $b^{3}df$ & \ldots & \\
$a^{2}$ $b^{2}c^{2}f$ & $+$ & $3$ & $a^{2}$ $b^{2}ce^{2}$ & $+$ & $4$ & $b^{4}cef$ & $+$ & $4$ & $c^{2}df^{2}$ & $-$ & $4$ \\
$a^{2}$ $c^{3}f$ & \ldots & & $a^{2}$ $c^{2}de$ & $-$ & $2$ & $c^{2}e^{3}$ & $-$ & $1$ & $cde^{2}f$ & $+$ & $2$ \\
$b^{3}d^{2}f$ & \ldots & & $b^{3}def$ & $-$ & $8$ & $c^{2}def$ & $-$ & $8$ & $d^{3}ef$ & $-$ & $12$ \\
$b^{2}c^{2}df$ & $+$ & $6$ & $b^{2}cd^{2}e$ & $+$ & $12$ & $cd^{3}f$ & $+$ & $2$ & $d^{4}e$ & $+$ & $6$ \\
$bc^{3}ef$ & $+$ & $4$ & $bc^{2}d^{2}f$ & $-$ & $6$ & $c^{3}e^{2}$ & $+$ & $6$ & $c^{2}d^{2}e^{2}$ & $-$ & $4$ \\
\hline
\end{tabular}
\end{center}
\end{footnotesize}

\hfill $(a,b,c,d,e,f\between x,y)^{4}$

\vspace{2ex}

\begin{center}
\textbf{O.\quad No.\ 90.}
\end{center}

\begin{footnotesize}
\begin{center}
\begin{tabular}{|r@{\hspace{0.4em}}c@{\hspace{0.4em}}r|r@{\hspace{0.4em}}c@{\hspace{0.4em}}r|r@{\hspace{0.4em}}c@{\hspace{0.4em}}r|}
\hline
$a^{5}$ $b^{2}f^{2}$ & $+$ & $1$ & $a^{4}$ $b^{2}df$ & $-$ & $1$ & $a^{3}$ $b^{3}f^{2}$ & \ldots & \\
$a^{4}$ $c^{2}f^{2}$ & $+$ & $4$ & $a^{4}$ $cdef$ & $+$ & $4$ & $a^{3}$ $b^{2}ef$ & \ldots & \\
$a^{4}$ $d^{3}f$ & $-$ & $1$ & $a^{4}$ $d^{2}e^{2}$ & $-$ & $7$ & $a^{3}$ $bc^{2}f$ & $-$ & $3$ \\
$a^{3}$ $b^{2}d^{2}f$ & $-$ & $6$ & $a^{3}$ $b^{2}de^{2}$ & $-$ & $8$ & $a^{3}$ $bcd^{2}f$ & $+$ & $18$ \\
$a^{3}$ $c^{3}df$ & $+$ & $7$ & $a^{3}$ $c^{2}d^{2}e$ & $+$ & $2$ & $a^{2}$ $b^{4}f^{2}$ & \ldots & \\
$a^{2}$ $b^{2}c^{2}f$ & $+$ & $3$ & $a^{2}$ $b^{2}d^{3}$ & $+$ & $9$ & $a^{2}$ $bc^{3}e$ & $-$ & $8$ \\
\hline
\end{tabular}
\end{center}
\end{footnotesize}

\hfill $(a,b,c,d,e,f\between x,y)^{5}$

\clearpage
\setcounter{page}{285}

\begin{center}
\textbf{P.\quad No.\ 91.}
\end{center}

\begin{footnotesize}
\begin{center}
\begin{tabular}{|r@{\hspace{0.3em}}c@{\hspace{0.3em}}r|r@{\hspace{0.3em}}c@{\hspace{0.3em}}r|r@{\hspace{0.3em}}c@{\hspace{0.3em}}r|r@{\hspace{0.3em}}c@{\hspace{0.3em}}r|r@{\hspace{0.3em}}c@{\hspace{0.3em}}r|r@{\hspace{0.3em}}c@{\hspace{0.3em}}r|}
\hline
$a^{4}$ $b^{2}df^{3}$ & \ldots & & $a^{3}$ $b^{4}f^{3}$ & $-$ & $1$ & $a^{3}$ $b^{3}cf^{3}$ & $-$ & $1$ & $a^{2}$ $b^{5}f^{3}$ & \ldots & & $a^{2}$ $b^{4}df^{3}$ & \ldots & \\
$a^{4}$ $b^{2}e^{2}f^{2}$ & \ldots & & $a^{3}$ $b^{3}df^{2}$ & $-$ & $2$ & $a^{3}$ $b^{3}e^{2}f$ & $+$ & $6$ & $b^{6}f^{3}$ & \ldots & & $b^{5}df^{2}$ & \ldots & \\
$a^{4}$ $c^{3}f^{3}$ & $-$ & $1$ & $a^{3}$ $b^{2}c^{2}f^{2}$ & $-$ & $3$ & $a^{3}$ $b^{2}cdef$ & $+$ & $11$ & $a^{2}$ $b^{3}c^{2}f^{2}$ & $+$ & $5$ & $c^{4}f^{3}$ & $-$ & $1$ \\
$a^{4}$ $c^{2}df^{2}$ & $+$ & $5$ & $a^{3}$ $b^{2}ce^{3}$ & $-$ & $5$ & $a^{3}$ $b^{2}d^{2}f$ & $-$ & $8$ & $c^{3}df^{2}$ & $+$ & $5$ & $c^{3}e^{2}f$ & $-$ & $5$ \\
$a^{4}$ $cdef^{2}$ & $-$ & $8$ & $a^{3}$ $b^{2}d^{2}e^{2}$ & $+$ & $2$ & $a^{3}$ $b^{2}de^{3}$ & $+$ & $6$ & $c^{2}d^{2}f$ & $-$ & $6$ & $c^{2}de^{3}$ & $+$ & $4$ \\
$a^{4}$ $d^{4}f$ & $+$ & $6$ & $a^{3}$ $bc^{3}f^{2}$ & $+$ & $7$ & $a^{3}$ $bc^{2}df$ & $+$ & $30$ & $cd^{3}f$ & $+$ & $6$ & $d^{4}e$ & $-$ & $1$ \\
\hline
\end{tabular}
\end{center}
\end{footnotesize}

\hfill $(a,b,c,d,e,f\between x,y)^{7}$

\vspace{2ex}
\begin{center}
\textbf{Q.\quad No.\ 25.\quad Q$'$.\quad No.\ 26.}
\end{center}

\begin{footnotesize}
\begin{center}
\begin{tabular}{|r@{\hspace{0.4em}}c@{\hspace{0.4em}}r|r@{\hspace{0.4em}}c@{\hspace{0.4em}}r||r@{\hspace{0.4em}}c@{\hspace{0.4em}}r|r@{\hspace{0.4em}}c@{\hspace{0.4em}}r|}
\hline
$a^{4}$ $b^{2}f^{4}$ & \ldots & & $+$ & $1$ & $a^{5}$ $b^{2}e^{4}$ & $+$ & $27$ & $-$ & $3375$ \\
$a^{3}$ $b^{3}ef^{3}$ & \ldots & & $-$ & $20$ & $a^{4}$ $b^{2}cdf^{3}$ & $-$ & $48$ & $+$ & $5760$ \\
$b^{3}cdf^{3}$ & $+$ & $1$ & $-$ & $120$ & $c^{5}f$ & $+$ & $3$ & $-$ & $600$ \\
$c^{3}ef^{2}$ & $-$ & $1$ & $+$ & $160$ & $cd^{3}ef$ & $+$ & $106$ & $-$ & $16000$ \\
$d^{2}f^{2}$ & $-$ & $3$ & $+$ & $360$ & $cde^{3}$ & $-$ & $81$ & $+$ & $9000$ \\
$d^{2}e^{2}f$ & $+$ & $5$ & $-$ & $640$ & $d^{4}f$ & $-$ & $38$ & $+$ & $6400$ \\
$de^{4}$ & $-$ & $2$ & $+$ & $256$ & $d^{3}e^{2}$ & $+$ & $38$ & $-$ & $4000$ \\
\hline
\end{tabular}
\end{center}
\end{footnotesize}

\hfill $(a,b,c,d,e,f\between x,y)^{8}$

\clearpage
\setcounter{page}{286}

\begin{center}
\textbf{R.\quad No.\ 92.}
\end{center}

\begin{footnotesize}
\begin{center}
\begin{tabular}{|r@{\hspace{0.4em}}c@{\hspace{0.4em}}r|r@{\hspace{0.4em}}c@{\hspace{0.4em}}r|r@{\hspace{0.4em}}c@{\hspace{0.4em}}r|}
\hline
$a^{5}$ $b^{2}f^{2}$ & \ldots & & $a^{4}$ $b^{3}f^{3}$ & \ldots & & $a^{4}$ $b^{2}df^{2}$ & \ldots \\
$a^{5}$ $cef^{2}$ & \ldots & & $a^{4}$ $b^{2}ef^{2}$ & \ldots & & $a^{4}$ $bcd^{2}f$ & \ldots \\
$a^{5}$ $d^{2}f$ & \ldots & & $a^{4}$ $bc^{2}f^{2}$ & \ldots & & $a^{4}$ $c^{3}df$ & \ldots \\
$a^{4}$ $b^{2}d^{2}f$ & \ldots & & $a^{3}$ $b^{4}f^{2}$ & \ldots & & $a^{3}$ $b^{3}c^{2}f$ & \ldots \\
$a^{3}$ $b^{2}c^{2}df$ & \ldots & & $a^{3}$ $b^{2}d^{4}$ & \ldots & & $a^{2}$ $b^{5}f^{2}$ & \ldots \\
\hline
\end{tabular}
\end{center}
\end{footnotesize}

\hfill $(a,b,c,d,e,f\between x,y)^{2}$

\vspace{2ex}
\begin{center}
\textbf{S.\quad No.\ 93.\quad S$'$.\quad No.\ 93 bis.}
\end{center}

\begin{footnotesize}
\begin{center}
\begin{tabular}{|r@{\hspace{0.4em}}c@{\hspace{0.4em}}r|r@{\hspace{0.4em}}c@{\hspace{0.4em}}r||r@{\hspace{0.4em}}c@{\hspace{0.4em}}r|r@{\hspace{0.4em}}c@{\hspace{0.4em}}r|}
\hline
$a^{5}$ $b^{2}ef^{3}$ & \ldots & & \ldots & & $a^{5}$ $b^{2}d^{2}f^{2}$ & \ldots & & \ldots & \\
$a^{4}$ $b^{3}f^{3}$ & \ldots & & \ldots & & $a^{4}$ $b^{3}df^{2}$ & \ldots & & \ldots & \\
$a^{4}$ $bc^{2}f^{3}$ & \ldots & & \ldots & & $a^{4}$ $b^{2}c^{2}f^{2}$ & \ldots & & \ldots & \\
$a^{3}$ $b^{2}cdf^{2}$ & \ldots & & \ldots & & $a^{3}$ $b^{4}f^{2}$ & \ldots & & \ldots & \\
$a^{3}$ $c^{3}f^{2}$ & \ldots & & \ldots & & $a^{3}$ $b^{2}d^{3}f$ & \ldots & & \ldots & \\
\hline
\end{tabular}
\end{center}
\end{footnotesize}

\hfill $(a,b,c,d,e,f\between x,y)^{3}$

\vspace{2ex}
\begin{center}
\textbf{T.\quad No.\ 94.}
\end{center}

\begin{footnotesize}
\begin{center}
\begin{tabular}{|r@{\hspace{0.4em}}c@{\hspace{0.4em}}r|r@{\hspace{0.4em}}c@{\hspace{0.4em}}r||r@{\hspace{0.4em}}c@{\hspace{0.4em}}r|r@{\hspace{0.4em}}c@{\hspace{0.4em}}r|}
\hline
$a^{5}$ $b^{2}df^{2}$ & \ldots & & $-$ & $18$ & $a^{4}$ $b^{3}df^{3}$ & $-$ & $75$ & $a^{5}$ $b^{2}d^{2}ef$ & $-$ & $880$ \\
$c^{2}f^{4}$ & \ldots & & $+$ & $150$ & $c^{3}e^{2}f^{2}$ & $-$ & $3$ & $cd^{2}e^{3}$ & $+$ & $765$ \\
$c^{2}ef^{3}$ & $+$ & $1$ & $-$ & $72$ & $c^{3}def^{2}$ & $+$ & $108$ & $d^{4}ef$ & $+$ & $148$ \\
$cde^{2}f^{2}$ & $-$ & $2$ & $+$ & $198$ & $c^{3}de^{2}f$ & $-$ & $308$ & $d^{3}e^{3}$ & $+$ & $175$ \\
$ce^{3}f$ & $+$ & $1$ & $-$ & $315$ & $c^{3}e^{3}$ & $+$ & $112$ & $b^{3}cf^{3}$ & $-$ & $24$ \\
$d^{3}f^{2}$ & $-$ & $3$ & $+$ & $129$ & $c^{2}d^{2}f^{2}$ & $-$ & $663$ & $c^{4}ef^{2}$ & $-$ & $513$ \\
$d^{2}e^{2}f$ & $+$ & $8$ & $-$ & $783$ & $c^{2}d^{2}ef$ & $+$ & $283$ & $c^{3}d^{2}f^{2}$ & $-$ & $283$ \\
$de^{4}$ & $-$ & $7$ & $+$ & $9$ & $c^{2}de^{3}$ & $-$ & $914$ & $c^{3}de^{2}f$ & $+$ & $1986$ \\
$e^{5}$ & $+$ & $2$ & $+$ & $48$ & $c^{2}d^{2}e^{2}$ & $+$ & $280$ & $c^{2}d^{4}$ & $-$ & $1365$ \\
\hline
\end{tabular}
\end{center}
\end{footnotesize}

\hfill $(a,b,c,d,e,f\between x,y)^{1}$

\clearpage
\setcounter{page}{287}

\begin{center}
\textbf{U.\quad No.\ 29.}
\end{center}

\begin{footnotesize}
\begin{center}
\begin{tabular}{|r@{\hspace{0.4em}}c@{\hspace{0.4em}}r|r@{\hspace{0.4em}}c@{\hspace{0.4em}}r|r@{\hspace{0.4em}}c@{\hspace{0.4em}}r|}
\hline
$a^{6}$ $c^{3}f^{3}$ & \ldots & & $a^{5}$ $b^{2}df^{3}$ & \ldots & & $a^{4}$ $b^{4}f^{3}$ & \ldots \\
$a^{6}$ $c^{2}df^{2}$ & \ldots & & $a^{5}$ $b^{2}e^{2}f^{2}$ & \ldots & & $a^{4}$ $b^{3}cdf^{2}$ & \ldots \\
$a^{6}$ $cd^{2}f$ & \ldots & & $a^{5}$ $bc^{2}f^{3}$ & \ldots & & $a^{4}$ $b^{2}c^{2}df$ & \ldots \\
$a^{6}$ $d^{3}$ & \ldots & & $a^{5}$ $bcdef^{2}$ & \ldots & & $a^{4}$ $b^{2}cd^{2}f$ & \ldots \\
$a^{5}$ $b^{2}f^{3}$ & \ldots & & $a^{5}$ $bd^{2}ef$ & \ldots & & $a^{4}$ $b^{2}de^{3}$ & \ldots \\
\hline
\end{tabular}
\end{center}
\end{footnotesize}

\hfill $(a,b,c,d,e,f\between x,y)^{0}$ Invariant.

\vspace{2ex}
\begin{center}
\textbf{V.\quad No.\ 95.}
\end{center}

\begin{footnotesize}
\begin{center}
\begin{tabular}{|r@{\hspace{0.4em}}c@{\hspace{0.4em}}r|r@{\hspace{0.4em}}c@{\hspace{0.4em}}r||r@{\hspace{0.4em}}c@{\hspace{0.4em}}r|}
\hline
$a^{5}$ $b^{2}d^{2}f^{3}$ & $+$ & $1620$ & $a^{4}$ $b^{4}df^{3}$ & $-$ & $276$ & $a^{5}$ $b^{2}c^{2}de^{4}$ & $+$ & $7100$ \\
$c^{4}de^{2}f^{2}$ & $+$ & $3408$ & $d^{2}e^{3}f^{3}$ & $+$ & $908$ & $c^{2}d^{5}ef$ & $+$ & $12440$ \\
$c^{4}e^{3}f$ & $+$ & $2156$ & $d^{2}e^{4}f^{2}$ & $-$ & $810$ & $c^{2}d^{4}e^{2}f$ & $-$ & $5200$ \\
$c^{3}d^{2}ef^{2}$ & $-$ & $15060$ & $d^{6}$ & \ldots & & $c^{2}d^{3}e^{3}$ & $-$ & $5340$ \\
$c^{3}d^{2}e^{2}f$ & $-$ & $2360$ & $b^{3}c^{3}f^{4}$ & $-$ & $12$ & $c^{2}d^{2}e^{4}$ & $-$ & $11900$ \\
$c^{3}de^{4}$ & $+$ & $9260$ & $c^{5}def^{3}$ & $+$ & $832$ & $c^{3}d^{3}e^{2}$ & $+$ & $10800$ \\
$c^{2}d^{3}f^{2}$ & $+$ & $4336$ & $c^{5}e^{2}f^{2}$ & $+$ & $1244$ & $c^{3}d^{2}e^{3}$ & $-$ & $2250$ \\
$c^{2}d^{2}e^{2}f$ & $+$ & $15220$ & $c^{4}d^{2}f^{3}$ & $+$ & $1112$ & $b^{3}c^{3}ef^{2}$ & $+$ & $486$ \\
$c^{2}de^{4}$ & $+$ & $19920$ & $c^{4}de^{2}f^{2}$ & $-$ & $168$ & $c^{4}d^{2}f^{2}$ & $+$ & $810$ \\
$c^{2}d^{3}ef$ & $-$ & $5808$ & $c^{4}d^{2}ef$ & $-$ & $3510$ & $c^{4}def^{2}$ & $+$ & $3330$ \\
$c^{2}d^{2}e^{3}$ & $-$ & $22740$ & $c^{4}de^{3}$ & $-$ & $1350$ & $c^{4}e^{3}$ & $-$ & $750$ \\
$cd^{4}f$ & $-$ & $90$ & $cd^{5}f^{2}$ & $+$ & $1172$ & $c^{3}d^{2}ef$ & $-$ & $8160$ \\
$cd^{3}e^{2}$ & $+$ & $10080$ & $d^{6}ef$ & $-$ & $2060$ & $c^{5}d^{2}e^{2}$ & $+$ & $5400$ \\
$d^{5}e$ & $-$ & $1350$ & $d^{4}e^{3}$ & $+$ & $450$ & $c^{5}de^{3}$ & $+$ & $3100$ \\
$b^{3}c^{2}df^{3}$ & $-$ & $864$ & $b^{3}c^{2}ef^{3}$ & $-$ & $240$ & $c^{6}d^{2}e^{2}$ & $+$ & $13900$ \\
$c^{5}f^{3}$ & $-$ & $1008$ & $c^{4}df^{2}$ & $-$ & $1620$ & $c^{6}de^{2}$ & $-$ & $9100$ \\
$c^{4}ef^{2}$ & $+$ & $6624$ & $c^{4}de^{2}f^{2}$ & $-$ & $3480$ & $c^{6}e^{3}$ & $+$ & $1800$ \\
$c^{4}e^{2}f$ & $-$ & $5344$ & $c^{4}e^{3}f$ & $-$ & $430$ & $b^{3}c^{4}df^{2}$ & $+$ & $162$ \\
$c^{3}d^{2}e^{2}f$ & $+$ & $2952$ & $c^{3}d^{3}f^{2}$ & $+$ & $876$ & $c^{8}f^{2}$ & $-$ & $810$ \\
$c^{3}de^{4}$ & $-$ & $3440$ & $c^{3}d^{2}ef^{2}$ & $+$ & $1224$ & $c^{7}def$ & $+$ & $1620$ \\
$c^{3}e^{5}$ & $+$ & $1720$ & $c^{3}def^{2}$ & $-$ & $68$ & $c^{7}e^{2}$ & $+$ & $1500$ \\
$c^{2}d^{4}f$ & $+$ & $4920$ & $c^{2}d^{2}e^{2}f$ & $-$ & $12$ & $c^{6}d^{2}f$ & $-$ & $600$ \\
$c^{2}d^{3}e^{2}$ & $-$ & $4768$ & $c^{2}d^{2}e^{3}$ & $-$ & $2800$ & $c^{6}de^{2}$ & $-$ & $3150$ \\
$c^{2}d^{2}e^{3}$ & $-$ & $16520$ & $c^{2}de^{4}$ & $+$ & $3840$ & $c^{6}e^{3}$ & $+$ & $2000$ \\
$c^{2}df^{4}$ & $+$ & $1920$ & $cd^{4}f^{2}$ & $-$ & $8540$ & $c^{5}d^{3}$ & $-$ & $400$ \\
$c^{2}d^{2}e^{2}$ & $+$ & $19440$ & $c^{2}d^{3}e^{2}$ & $-$ & $10100$ & \multicolumn{3}{c|}{} \\
$cd^{5}$ & $-$ & $9540$ & $c^{2}d^{2}e^{3}f$ & $-$ & $6240$ & \multicolumn{3}{c|}{} \\
$d^{6}$ & $+$ & $1620$ & $c^{2}de^{4}$ & $+$ & $23300$ & \multicolumn{3}{c|}{} \\
$b^{3}cf^{4}$ & $-$ & $162$ & $c^{2}d^{4}f$ & $+$ & $3640$ & \multicolumn{3}{c|}{} \\
$b^{2}c^{2}ef^{3}$ & $+$ & $918$ & $c^{2}d^{3}e^{2}$ & $+$ & $8800$ & \multicolumn{3}{c|}{} \\
$c^{6}f^{2}$ & $-$ & $1956$ & $c^{2}d^{2}e^{3}$ & $-$ & $750$ & \multicolumn{3}{c|}{} \\
$c^{5}df^{2}$ & $+$ & $240$ & $d^{5}e$ & $+$ & $450$ & \multicolumn{3}{c|}{} \\
\hline
\multicolumn{3}{|c|}{$\pm\,6$} & \multicolumn{3}{c|}{$\pm\,169$\quad$\pm\,625$} & \multicolumn{3}{c|}{$\pm\,424$\quad$\pm\,1124$} \\
\hline
\end{tabular}
\end{center}
\end{footnotesize}

\hfill $(a,b,c,d,e,f\between x,y)^{1}$

\clearpage
\setcounter{page}{288}

\begin{center}
\textbf{W.\quad No.\ 29A.}
\end{center}

\begin{footnotesize}
\begin{center}
\begin{tabular}{|r@{\hspace{0.4em}}c@{\hspace{0.4em}}r|r@{\hspace{0.4em}}c@{\hspace{0.4em}}r|r@{\hspace{0.4em}}c@{\hspace{0.4em}}r|}
\hline
$a^{9}$ $c^{4}f^{4}$ & \ldots & & $a^{8}$ $b^{2}d^{2}f^{4}$ & \ldots & & $a^{7}$ $b^{4}f^{4}$ & \ldots \\
$a^{8}$ $c^{3}df^{3}$ & \ldots & & $a^{7}$ $b^{3}cdf^{3}$ & \ldots & & $a^{6}$ $b^{6}f^{4}$ & \ldots \\
$a^{7}$ $c^{2}d^{2}f^{2}$ & \ldots & & $a^{6}$ $b^{4}c^{2}f^{2}$ & \ldots & & $a^{5}$ $b^{6}df^{2}$ & \ldots \\
$a^{6}$ $cd^{3}f$ & \ldots & & $a^{5}$ $b^{5}d^{2}f$ & \ldots & & $a^{4}$ $b^{8}f^{2}$ & \ldots \\
$a^{5}$ $d^{4}$ & \ldots & & $a^{4}$ $b^{6}cd^{2}$ & \ldots & & $a^{3}$ $b^{8}e^{2}$ & \ldots \\
\hline
\end{tabular}
\end{center}
\end{footnotesize}

\hfill $(a,b,c,d,e,f\between x,y)^{0}$ Invariant.

\vspace{2ex}

The covariants Q and Q$'$ (Nos.\ 25 and 26) are given in the Second Memoir, 141, but are there presented in a different form. The covariant S$'$ (No.\ 93 bis) is an alternative form of S; and similarly Q$'$ is an alternative form of Q. The remaining covariants R, T, U, V, W are here presented with their full series of literal terms; the tables exhibit the literal terms with their numerical coefficients, arranged according to the powers of $a$ and $b$, as explained above.

The numerical verifications at the foot of each table show the sums of the positive and negative coefficients for each set of terms containing the same power of $a$; and in the last line, the sum for the entire covariant. Where these sums are equal (with opposite signs), they are marked with the sign $\pm$; where they are unequal, the actual values are given.

\clearpage
\setcounter{page}{289}

The cases of inequality in the numerical verifications are as follows:

\begin{center}
\begin{tabular}{|l|c|c|l|}
\hline
\multicolumn{1}{|c|}{Covariant} & Positive sum & Negative sum & \multicolumn{1}{c|}{Difference} \\
\hline
P (No.\ 91) & $128505$ & $128505$ & $0$ \\
Q (No.\ 25) & $776$ & $780$ & $-4$ \\
Q$'$ (No.\ 26) & $21256$ & $21250$ & $+6$ \\
R (No.\ 92) & $68056$ & $68660$ & $-4$ \\
S (No.\ 93) & $37816$ & $37815$ & $+1$ \\
T (No.\ 94) & $20196$ & $20196$ & $0$ \\
U (No.\ 29) & $128505$ & $128505$ & $0$ \\
V (No.\ 95) & varies & varies & see table \\
W (No.\ 29A) & $1124$ & $1124$ & $0$ \\
\hline
\end{tabular}
\end{center}

The verification for the entire covariant P is $\pm\,128505$; and the several verifications for Q are $1=1$, $776-780=-4$, $21256-21250=+6$, $68056-68060=-4$, $37816-37815=+1$, giving in sum $128505-128505=0$.

\clearpage
\setcounter{page}{290}

\begin{center}
\textbf{Note on the Covariants of the Binary Quintic.}
\end{center}

The foregoing tables complete the series of covariants of the binary quintic, as given in my several Memoirs on Quantics. The entire system consists of 23 covariants (including the quintic itself and the four invariants G, Q, U, W); and these are connected by a system of syzygies, which I have investigated in the Tenth Memoir on Quantics. The present arrangement exhibits the covariants in a compendious form, and facilitates the verification of the numerical coefficients.

I have thought it desirable to preserve the original notation, in which the covariants are denoted by the capital letters A, B, C, \ldots W; but I have added references to the Memoirs in which each covariant was first obtained, and the number of the table in the original publication. The alternative forms Q$'$ and S$'$ are new; they were obtained by substituting for the original forms the expressions which arise from the syzygies connecting the covariants.

The numerical tables, 142, give the number of terms of each degree and weight for the several quantics; and by means of these tables, the number of terms in each covariant can be ascertained beforehand. This affords a complete verification of the literal terms; and the numerical verification of the coefficients is effected by the sums of the positive and negative coefficients, as shown in the tables.

\clearpage
\setcounter{page}{291}

The syzygies which connect the covariants of the binary quintic are of great importance in the theory. They show that certain combinations of the covariants are identically zero; and they thus establish relations between the coefficients, which are of use in the reduction of the quintic to its canonical form. The principal syzygies are those which connect the covariants of the same degree and order; and these have been fully investigated in the Tenth Memoir.

I have also given in the present reprint the numerical tables supplementary to the Second Memoir, 142; these tables exhibit the number of literal terms of each degree and weight for the cubic, quartic, quintic, sextic, septimic, and octavic functions. The tables are carried to a sufficient extent to include the highest invariants of each function; and they serve as a verification of the completeness of the covariant tables.

The theory of the covariants of the binary quintic is thus brought to a satisfactory completion. The entire system of 23 covariants has been obtained and verified; the syzygies connecting them have been investigated; and the numerical tables afford a ready means of ascertaining the number of terms in any covariant of a given degree and order. The results are presented in a form which admits of easy reference; and it is hoped that they will be of service in future investigations on this important subject.

\clearpage
\setcounter{page}{292}

\begin{center}
\textbf{Additional Note.}
\end{center}

Since the above was written, I have found it convenient to make some further changes in the notation and arrangement of the covariants. In particular, I have adopted a uniform system of lettering, in which the covariants are denoted by Roman capitals, without suffixes; and I have arranged them in the order of their degrees. This facilitates the reference to any particular covariant; and it renders the notation consistent with that employed in my later Memoirs on Quantics.

I have also made some corrections in the numerical coefficients, which have been detected in the process of verification. These corrections are comparatively few in number; and they do not affect the general theory. The corrected values are given in the present reprint; and the original values are indicated in the notes.

The present Memoir completes the series of investigations on the binary quintic, which I have published in the Philosophical Transactions. The results are collected in a form which admits of easy reference; and it is believed that they are substantially correct. The theory of the covariants of the binary quintic is a necessary preliminary to the theory of the equations of the fifth degree; and it is hoped that the present investigation will contribute to the advancement of this important branch of algebra.

\clearpage
\setcounter{page}{293}

\begin{center}
\textbf{T.\quad No.\ 94 (continued).}
\end{center}

\begin{footnotesize}
\begin{center}
\begin{tabular}{|r@{\hspace{0.3em}}c@{\hspace{0.3em}}r|r@{\hspace{0.3em}}c@{\hspace{0.3em}}r|r@{\hspace{0.3em}}c@{\hspace{0.3em}}r|r@{\hspace{0.3em}}c@{\hspace{0.3em}}r|}
\hline
$a^{5}$ $b^{2}d^{2}f^{3}$ & $+$ & $1620$ & $a^{4}$ $b^{3}cd^{2}f^{3}$ & $-$ & $18$ & $a^{3}$ $b^{4}df^{3}$ & $-$ & $75$ & $a^{2}$ $b^{5}d^{2}ef$ & $-$ & $880$ \\
$c^{2}de^{2}f^{2}$ & $+$ & $3408$ & $c^{2}e^{3}f^{2}$ & $+$ & $150$ & $c^{3}e^{2}f^{2}$ & $-$ & $3$ & $c^{2}d^{2}e^{3}$ & $+$ & $765$ \\
$c^{2}e^{4}f$ & $+$ & $2156$ & $c^{2}e^{2}f^{3}$ & $-$ & $72$ & $c^{3}def^{2}$ & $+$ & $108$ & $cd^{2}e^{4}$ & $+$ & $148$ \\
$c^{2}d^{3}f^{2}$ & $-$ & $15060$ & $def^{4}$ & $+$ & $198$ & $c^{3}de^{2}f$ & $-$ & $308$ & $d^{3}e^{3}$ & $+$ & $175$ \\
$c^{2}d^{2}e^{2}f$ & $-$ & $2360$ & $d^{2}e^{2}f^{2}$ & $-$ & $315$ & $c^{3}e^{3}$ & $+$ & $112$ & $b^{3}cf^{3}$ & $-$ & $24$ \\
$c^{2}de^{4}$ & $+$ & $9260$ & $d^{2}e^{3}f$ & $+$ & $129$ & $c^{2}d^{2}f^{2}$ & $-$ & $663$ & $c^{4}ef^{2}$ & $-$ & $513$ \\
$c^{2}d^{4}f$ & $+$ & $4336$ & $de^{5}$ & $-$ & $6$ & $c^{2}def$ & $-$ & $783$ & $c^{3}d^{2}f^{2}$ & $-$ & $283$ \\
$c^{2}d^{3}e^{2}$ & $+$ & $15220$ & $c^{3}f^{4}$ & $-$ & $1$ & $c^{2}de^{3}$ & $+$ & $9$ & $c^{3}de^{2}f$ & $+$ & $1986$ \\
$c^{2}d^{2}e^{4}$ & $+$ & $19920$ & $c^{2}def^{3}$ & $+$ & $66$ & $c^{2}d^{2}e^{2}$ & $+$ & $48$ & $c^{2}d^{4}$ & $-$ & $280$ \\
$c^{2}d^{3}ef$ & $-$ & $5808$ & $c^{2}e^{2}f^{2}$ & $-$ & $114$ & $cd^{3}f$ & $+$ & $216$ & $c^{2}d^{3}e^{2}$ & $-$ & $1365$ \\
$c^{2}d^{2}e^{3}$ & $-$ & $22740$ & $cdef^{3}$ & $-$ & $153$ & $cd^{2}e^{2}$ & $-$ & $252$ & $c^{2}d^{2}e^{3}$ & $+$ & $527$ \\
$cd^{5}$ & $-$ & $90$ & $ce^{3}f$ & $+$ & $108$ & $b^{3}c^{2}f^{3}$ & $+$ & $27$ & $c^{2}de^{4}$ & $-$ & $340$ \\
$cd^{4}e^{2}$ & $+$ & $10080$ & $d^{2}f^{3}$ & $-$ & $81$ & $c^{4}f^{3}$ & $-$ & $1$ & $cd^{4}f$ & $-$ & $700$ \\
$d^{6}e$ & $-$ & $1350$ & $d^{2}ef^{2}$ & $+$ & $63$ & $c^{3}df^{2}$ & $+$ & $294$ & $d^{5}$ & $+$ & $60$ \\
$b^{3}c^{2}df^{3}$ & $-$ & $864$ & $de^{3}f$ & $+$ & $28$ & $c^{3}ef^{2}$ & $-$ & $208$ & $b^{3}c^{2}ef^{2}$ & $+$ & $153$ \\
$c^{4}f^{3}$ & $-$ & $1008$ & $e^{5}$ & $-$ & $3$ & $c^{3}de^{2}f$ & $-$ & $93$ & $c^{4}d^{2}f^{2}$ & $-$ & $1032$ \\
$c^{4}ef^{2}$ & $+$ & $6624$ & $b^{2}c^{2}f^{4}$ & $+$ & $1$ & $c^{2}d^{2}f^{2}$ & $-$ & $570$ & $c^{4}de^{2}f$ & $-$ & $1098$ \\
$c^{4}e^{2}f$ & $-$ & $5344$ & $b^{2}cdef^{3}$ & $+$ & $5$ & $c^{2}d^{2}ef$ & $+$ & $32$ & $c^{4}e^{3}$ & $-$ & $40$ \\
$c^{3}d^{2}f^{2}$ & $+$ & $2952$ & $b^{2}ce^{2}f^{2}$ & $-$ & $8$ & $c^{2}de^{3}$ & $-$ & $42$ & $c^{3}d^{2}ef$ & $-$ & $1662$ \\
$c^{3}de^{3}$ & $-$ & $3440$ & $b^{2}d^{2}f^{3}$ & $+$ & $6$ & $cd^{3}f$ & $+$ & $1116$ & $c^{3}de^{3}$ & $+$ & $1025$ \\
$c^{3}e^{4}$ & $+$ & $1720$ & $b^{2}d^{2}ef^{2}$ & $-$ & $4$ & $cd^{2}e^{2}$ & $+$ & $224$ & $c^{2}d^{3}f$ & $+$ & $730$ \\
$c^{2}d^{4}f$ & $+$ & $4920$ & $b^{2}de^{2}f$ & $-$ & $18$ & $c^{2}d^{3}$ & $-$ & $798$ & $c^{2}d^{2}e^{2}$ & $+$ & $370$ \\
$c^{2}d^{3}e^{2}$ & $-$ & $4768$ & $bd^{3}f^{2}$ & $-$ & $9$ & $c^{2}d^{2}e^{2}$ & $+$ & $369$ & $c^{2}de^{3}$ & $-$ & $488$ \\
$c^{2}d^{2}e^{4}$ & $-$ & $16520$ & $bcdc^{2}f$ & $-$ & $12$ & $c^{2}de^{3}$ & $-$ & $486$ & $cd^{3}e$ & $-$ & $880$ \\
$c^{2}df^{4}$ & $+$ & $1920$ & $c^{4}f^{2}$ & $-$ & $8$ & $cd^{4}$ & $+$ & $81$ & $c^{2}d^{4}$ & $-$ & $735$ \\
$c^{2}d^{2}e^{2}$ & $+$ & $19440$ & $c^{3}df$ & $-$ & $5$ & $b^{2}c^{2}ef^{2}$ & $+$ & $5$ & $c^{2}d^{3}e$ & $-$ & $150$ \\
$cd^{5}$ & $-$ & $9540$ & $c^{3}e^{2}$ & $+$ & $7$ & $b^{2}cd^{2}f^{2}$ & $-$ & $6$ & $b^{2}c^{2}d^{2}f$ & $+$ & $81$ \\
$d^{6}$ & $+$ & $1620$ & $c^{2}d^{2}f$ & $-$ & $30$ & $b^{2}cde^{2}f$ & $-$ & $4$ & $b^{2}c^{2}de^{2}$ & $+$ & $243$ \\
$b^{3}cf^{4}$ & $-$ & $162$ & $c^{2}de^{2}$ & $+$ & $18$ & $b^{2}cd^{2}e^{2}$ & $+$ & $6$ & $b^{2}c^{3}e^{2}$ & $-$ & $30$ \\
$b^{2}c^{2}ef^{3}$ & $+$ & $918$ & $cd^{3}$ & $+$ & $8$ & $b^{2}d^{3}f$ & $-$ & $4$ & $b^{2}c^{2}d^{3}$ & $+$ & $102$ \\
$c^{6}f^{2}$ & $-$ & $1956$ & $bcdc^{3}$ & $-$ & $1$ & $b^{2}d^{3}e$ & $+$ & $8$ & $b^{2}d^{4}$ & $-$ & $18$ \\
$c^{5}df^{2}$ & $+$ & $240$ & $bc^{2}d^{2}$ & $-$ & $4$ & $bc^{3}df$ & $-$ & $6$ & $bc^{3}d^{2}e$ & $-$ & $12$ \\
\hline
\end{tabular}
\end{center}
\end{footnotesize}

\hfill $(a,b,c,d,e,f\between x,y)^{1}$

\clearpage
\setcounter{page}{294}

\begin{center}
\textbf{V.\quad No.\ 95 (continued).}
\end{center}

\begin{footnotesize}
\begin{center}
\begin{tabular}{|r@{\hspace{0.3em}}c@{\hspace{0.3em}}r|r@{\hspace{0.3em}}c@{\hspace{0.3em}}r|r@{\hspace{0.3em}}c@{\hspace{0.3em}}r|}
\hline
$a^{5}$ $b^{2}c^{2}d^{2}f^{3}$ & $+$ & $1620$ & $a^{4}$ $b^{4}df^{3}$ & $-$ & $276$ & $a^{5}$ $b^{2}c^{2}de^{4}$ & $+$ & $7100$ \\
$c^{4}de^{2}f^{2}$ & $+$ & $3408$ & $d^{2}e^{3}f^{3}$ & $+$ & $908$ & $c^{2}d^{5}ef$ & $+$ & $12440$ \\
$c^{4}e^{3}f$ & $+$ & $2156$ & $d^{2}e^{4}f^{2}$ & $-$ & $810$ & $c^{2}d^{4}e^{2}f$ & $-$ & $5200$ \\
$c^{3}d^{2}ef^{2}$ & $-$ & $15060$ & $d^{6}$ & \ldots & & $c^{2}d^{3}e^{3}$ & $-$ & $5340$ \\
$c^{3}d^{2}e^{2}f$ & $-$ & $2360$ & $b^{3}c^{3}f^{4}$ & $-$ & $12$ & $c^{2}d^{2}e^{4}$ & $-$ & $11900$ \\
$c^{3}de^{4}$ & $+$ & $9260$ & $c^{5}def^{3}$ & $+$ & $832$ & $c^{3}d^{3}e^{2}$ & $+$ & $10800$ \\
$c^{2}d^{3}f^{2}$ & $+$ & $4336$ & $c^{5}e^{2}f^{2}$ & $+$ & $1244$ & $c^{3}d^{2}e^{3}$ & $-$ & $2250$ \\
$c^{2}d^{2}e^{2}f$ & $+$ & $15220$ & $c^{4}d^{2}f^{3}$ & $+$ & $1112$ & $b^{3}c^{3}ef^{2}$ & $+$ & $486$ \\
$c^{2}de^{4}$ & $+$ & $19920$ & $c^{4}de^{2}f^{2}$ & $-$ & $168$ & $c^{4}d^{2}f^{2}$ & $+$ & $810$ \\
$c^{2}d^{3}ef$ & $-$ & $5808$ & $c^{4}d^{2}ef$ & $-$ & $3510$ & $c^{4}def^{2}$ & $+$ & $3330$ \\
$c^{2}d^{2}e^{3}$ & $-$ & $22740$ & $c^{4}de^{3}$ & $-$ & $1350$ & $c^{4}e^{3}$ & $-$ & $750$ \\
$cd^{4}f$ & $-$ & $90$ & $cd^{5}f^{2}$ & $+$ & $1172$ & $c^{3}d^{2}ef$ & $-$ & $8160$ \\
$cd^{3}e^{2}$ & $+$ & $10080$ & $d^{6}ef$ & $-$ & $2060$ & $c^{5}d^{2}e^{2}$ & $+$ & $5400$ \\
$d^{5}e$ & $-$ & $1350$ & $d^{4}e^{3}$ & $+$ & $450$ & $c^{5}de^{3}$ & $+$ & $3100$ \\
$b^{3}c^{2}df^{3}$ & $-$ & $864$ & $b^{3}c^{2}ef^{3}$ & $-$ & $240$ & $c^{6}d^{2}e^{2}$ & $+$ & $13900$ \\
$c^{5}f^{3}$ & $-$ & $1008$ & $c^{4}df^{2}$ & $-$ & $1620$ & $c^{6}de^{2}$ & $-$ & $9100$ \\
$c^{4}ef^{2}$ & $+$ & $6624$ & $c^{4}de^{2}f^{2}$ & $-$ & $3480$ & $c^{6}e^{3}$ & $+$ & $1800$ \\
$c^{4}e^{2}f$ & $-$ & $5344$ & $c^{4}e^{3}f$ & $-$ & $430$ & $b^{3}c^{4}df^{2}$ & $+$ & $162$ \\
$c^{3}d^{2}e^{2}f$ & $+$ & $2952$ & $c^{3}d^{3}f^{2}$ & $+$ & $876$ & $c^{8}f^{2}$ & $-$ & $810$ \\
$c^{3}de^{4}$ & $-$ & $3440$ & $c^{3}d^{2}ef^{2}$ & $+$ & $1224$ & $c^{7}def$ & $+$ & $1620$ \\
$c^{3}e^{5}$ & $+$ & $1720$ & $c^{3}def^{2}$ & $-$ & $68$ & $c^{7}e^{2}$ & $+$ & $1500$ \\
$c^{2}d^{4}f$ & $+$ & $4920$ & $c^{2}d^{2}e^{2}f$ & $-$ & $12$ & $c^{6}d^{2}f$ & $-$ & $600$ \\
$c^{2}d^{3}e^{2}$ & $-$ & $4768$ & $c^{2}d^{2}e^{3}$ & $-$ & $2800$ & $c^{6}de^{2}$ & $-$ & $3150$ \\
$c^{2}d^{2}e^{4}$ & $-$ & $16520$ & $c^{2}de^{4}$ & $+$ & $3840$ & $c^{6}e^{3}$ & $+$ & $2000$ \\
$c^{2}df^{4}$ & $+$ & $1920$ & $cd^{4}f^{2}$ & $-$ & $8540$ & $c^{5}d^{3}$ & $-$ & $400$ \\
$c^{2}d^{2}e^{2}$ & $+$ & $19440$ & $c^{2}d^{3}e^{2}$ & $-$ & $10100$ & \multicolumn{3}{c|}{} \\
$cd^{5}$ & $-$ & $9540$ & $c^{2}d^{2}e^{3}f$ & $-$ & $6240$ & \multicolumn{3}{c|}{} \\
$d^{6}$ & $+$ & $1620$ & $c^{2}de^{4}$ & $+$ & $23300$ & \multicolumn{3}{c|}{} \\
\hline
\multicolumn{3}{|c|}{$\pm\,12$} & \multicolumn{3}{c|}{$\pm\,395$\quad$\pm\,1650$} & \multicolumn{3}{c|}{$\pm\,6511$\quad$\pm\,11628$} \\
\hline
\multicolumn{9}{|c|}{$\pm\,20196$} \\
\hline
\end{tabular}
\end{center}
\end{footnotesize}

\hfill $(a,b,c,d,e,f\between x,y)^{1}$

\clearpage
\setcounter{page}{295}

\begin{center}
\textbf{V.\quad No.\ 95 (continued).}
\end{center}

\begin{footnotesize}
\begin{center}
\begin{tabular}{|r@{\hspace{0.3em}}c@{\hspace{0.3em}}r|r@{\hspace{0.3em}}c@{\hspace{0.3em}}r|r@{\hspace{0.3em}}c@{\hspace{0.3em}}r|}
\hline
$a^{4}$ $b^{4}d^{2}f^{3}$ & $+$ & $1620$ & $a^{3}$ $b^{5}df^{3}$ & $-$ & $276$ & $a^{2}$ $b^{6}e^{4}$ & $+$ & $7100$ \\
$c^{4}de^{2}f^{2}$ & $+$ & $3408$ & $d^{2}e^{3}f^{3}$ & $+$ & $908$ & $c^{2}d^{5}ef$ & $+$ & $12440$ \\
$c^{4}e^{3}f$ & $+$ & $2156$ & $d^{2}e^{4}f^{2}$ & $-$ & $810$ & $c^{2}d^{4}e^{2}f$ & $-$ & $5200$ \\
$c^{3}d^{2}ef^{2}$ & $-$ & $15060$ & $d^{6}$ & \ldots & & $c^{2}d^{3}e^{3}$ & $-$ & $5340$ \\
$c^{3}d^{2}e^{2}f$ & $-$ & $2360$ & $b^{3}c^{3}f^{4}$ & $-$ & $12$ & $c^{2}d^{2}e^{4}$ & $-$ & $11900$ \\
$c^{3}de^{4}$ & $+$ & $9260$ & $c^{5}def^{3}$ & $+$ & $832$ & $c^{3}d^{3}e^{2}$ & $+$ & $10800$ \\
$c^{2}d^{3}f^{2}$ & $+$ & $4336$ & $c^{5}e^{2}f^{2}$ & $+$ & $1244$ & $c^{3}d^{2}e^{3}$ & $-$ & $2250$ \\
$c^{2}d^{2}e^{2}f$ & $+$ & $15220$ & $c^{4}d^{2}f^{3}$ & $+$ & $1112$ & $b^{3}c^{3}ef^{2}$ & $+$ & $486$ \\
$c^{2}de^{4}$ & $+$ & $19920$ & $c^{4}de^{2}f^{2}$ & $-$ & $168$ & $c^{4}d^{2}f^{2}$ & $+$ & $810$ \\
$c^{2}d^{3}ef$ & $-$ & $5808$ & $c^{4}d^{2}ef$ & $-$ & $3510$ & $c^{4}def^{2}$ & $+$ & $3330$ \\
$c^{2}d^{2}e^{3}$ & $-$ & $22740$ & $c^{4}de^{3}$ & $-$ & $1350$ & $c^{4}e^{3}$ & $-$ & $750$ \\
$cd^{4}f$ & $-$ & $90$ & $cd^{5}f^{2}$ & $+$ & $1172$ & $c^{3}d^{2}ef$ & $-$ & $8160$ \\
$cd^{3}e^{2}$ & $+$ & $10080$ & $d^{6}ef$ & $-$ & $2060$ & $c^{5}d^{2}e^{2}$ & $+$ & $5400$ \\
$d^{5}e$ & $-$ & $1350$ & $d^{4}e^{3}$ & $+$ & $450$ & $c^{5}de^{3}$ & $+$ & $3100$ \\
$b^{3}c^{2}df^{3}$ & $-$ & $864$ & $b^{3}c^{2}ef^{3}$ & $-$ & $240$ & $c^{6}d^{2}e^{2}$ & $+$ & $13900$ \\
$c^{5}f^{3}$ & $-$ & $1008$ & $c^{4}df^{2}$ & $-$ & $1620$ & $c^{6}de^{2}$ & $-$ & $9100$ \\
$c^{4}ef^{2}$ & $+$ & $6624$ & $c^{4}de^{2}f^{2}$ & $-$ & $3480$ & $c^{6}e^{3}$ & $+$ & $1800$ \\
$c^{4}e^{2}f$ & $-$ & $5344$ & $c^{4}e^{3}f$ & $-$ & $430$ & $b^{3}c^{4}df^{2}$ & $+$ & $162$ \\
$c^{3}d^{2}e^{2}f$ & $+$ & $2952$ & $c^{3}d^{3}f^{2}$ & $+$ & $876$ & $c^{8}f^{2}$ & $-$ & $810$ \\
$c^{3}de^{4}$ & $-$ & $3440$ & $c^{3}d^{2}ef^{2}$ & $+$ & $1224$ & $c^{7}def$ & $+$ & $1620$ \\
$c^{3}e^{5}$ & $+$ & $1720$ & $c^{3}def^{2}$ & $-$ & $68$ & $c^{7}e^{2}$ & $+$ & $1500$ \\
$c^{2}d^{4}f$ & $+$ & $4920$ & $c^{2}d^{2}e^{2}f$ & $-$ & $12$ & $c^{6}d^{2}f$ & $-$ & $600$ \\
$c^{2}d^{3}e^{2}$ & $-$ & $4768$ & $c^{2}d^{2}e^{3}$ & $-$ & $2800$ & $c^{6}de^{2}$ & $-$ & $3150$ \\
$c^{2}d^{2}e^{4}$ & $-$ & $16520$ & $c^{2}de^{4}$ & $+$ & $3840$ & $c^{6}e^{3}$ & $+$ & $2000$ \\
$c^{2}df^{4}$ & $+$ & $1920$ & $cd^{4}f^{2}$ & $-$ & $8540$ & $c^{5}d^{3}$ & $-$ & $400$ \\
\hline
\end{tabular}
\end{center}
\end{footnotesize}

\clearpage
\setcounter{page}{296}

\begin{center}
\textbf{V.\quad No.\ 95 (concluded).}
\end{center}

\begin{footnotesize}
\begin{center}
\begin{tabular}{|r@{\hspace{0.3em}}c@{\hspace{0.3em}}r|r@{\hspace{0.3em}}c@{\hspace{0.3em}}r|r@{\hspace{0.3em}}c@{\hspace{0.3em}}r|}
\hline
$a^{3}$ $b^{6}d^{2}f^{3}$ & $+$ & $1620$ & $a^{2}$ $b^{7}df^{3}$ & $-$ & $276$ & $a$ $b^{8}e^{4}$ & $+$ & $7100$ \\
$c^{4}de^{2}f^{2}$ & $+$ & $3408$ & $d^{2}e^{3}f^{3}$ & $+$ & $908$ & $c^{2}d^{5}ef$ & $+$ & $12440$ \\
$c^{4}e^{3}f$ & $+$ & $2156$ & $d^{2}e^{4}f^{2}$ & $-$ & $810$ & $c^{2}d^{4}e^{2}f$ & $-$ & $5200$ \\
$c^{3}d^{2}ef^{2}$ & $-$ & $15060$ & $d^{6}$ & \ldots & & $c^{2}d^{3}e^{3}$ & $-$ & $5340$ \\
$c^{3}d^{2}e^{2}f$ & $-$ & $2360$ & $b^{3}c^{3}f^{4}$ & $-$ & $12$ & $c^{2}d^{2}e^{4}$ & $-$ & $11900$ \\
$c^{3}de^{4}$ & $+$ & $9260$ & $c^{5}def^{3}$ & $+$ & $832$ & $c^{3}d^{3}e^{2}$ & $+$ & $10800$ \\
$c^{2}d^{3}f^{2}$ & $+$ & $4336$ & $c^{5}e^{2}f^{2}$ & $+$ & $1244$ & $c^{3}d^{2}e^{3}$ & $-$ & $2250$ \\
$c^{2}d^{2}e^{2}f$ & $+$ & $15220$ & $c^{4}d^{2}f^{3}$ & $+$ & $1112$ & $b^{3}c^{3}ef^{2}$ & $+$ & $486$ \\
$c^{2}de^{4}$ & $+$ & $19920$ & $c^{4}de^{2}f^{2}$ & $-$ & $168$ & $c^{4}d^{2}f^{2}$ & $+$ & $810$ \\
$c^{2}d^{3}ef$ & $-$ & $5808$ & $c^{4}d^{2}ef$ & $-$ & $3510$ & $c^{4}def^{2}$ & $+$ & $3330$ \\
$c^{2}d^{2}e^{3}$ & $-$ & $22740$ & $c^{4}de^{3}$ & $-$ & $1350$ & $c^{4}e^{3}$ & $-$ & $750$ \\
$cd^{4}f$ & $-$ & $90$ & $cd^{5}f^{2}$ & $+$ & $1172$ & $c^{3}d^{2}ef$ & $-$ & $8160$ \\
$cd^{3}e^{2}$ & $+$ & $10080$ & $d^{6}ef$ & $-$ & $2060$ & $c^{5}d^{2}e^{2}$ & $+$ & $5400$ \\
$d^{5}e$ & $-$ & $1350$ & $d^{4}e^{3}$ & $+$ & $450$ & $c^{5}de^{3}$ & $+$ & $3100$ \\
$b^{3}c^{2}df^{3}$ & $-$ & $864$ & $b^{3}c^{2}ef^{3}$ & $-$ & $240$ & $c^{6}d^{2}e^{2}$ & $+$ & $13900$ \\
$c^{5}f^{3}$ & $-$ & $1008$ & $c^{4}df^{2}$ & $-$ & $1620$ & $c^{6}de^{2}$ & $-$ & $9100$ \\
$c^{4}ef^{2}$ & $+$ & $6624$ & $c^{4}de^{2}f^{2}$ & $-$ & $3480$ & $c^{6}e^{3}$ & $+$ & $1800$ \\
$c^{4}e^{2}f$ & $-$ & $5344$ & $c^{4}e^{3}f$ & $-$ & $430$ & $b^{3}c^{4}df^{2}$ & $+$ & $162$ \\
$c^{3}d^{2}e^{2}f$ & $+$ & $2952$ & $c^{3}d^{3}f^{2}$ & $+$ & $876$ & $c^{8}f^{2}$ & $-$ & $810$ \\
$c^{3}de^{4}$ & $-$ & $3440$ & $c^{3}d^{2}ef^{2}$ & $+$ & $1224$ & $c^{7}def$ & $+$ & $1620$ \\
$c^{3}e^{5}$ & $+$ & $1720$ & $c^{3}def^{2}$ & $-$ & $68$ & $c^{7}e^{2}$ & $+$ & $1500$ \\
$c^{2}d^{4}f$ & $+$ & $4920$ & $c^{2}d^{2}e^{2}f$ & $-$ & $12$ & $c^{6}d^{2}f$ & $-$ & $600$ \\
$c^{2}d^{3}e^{2}$ & $-$ & $4768$ & $c^{2}d^{2}e^{3}$ & $-$ & $2800$ & $c^{6}de^{2}$ & $-$ & $3150$ \\
$c^{2}d^{2}e^{4}$ & $-$ & $16520$ & $c^{2}de^{4}$ & $+$ & $3840$ & $c^{6}e^{3}$ & $+$ & $2000$ \\
$c^{2}df^{4}$ & $+$ & $1920$ & $cd^{4}f^{2}$ & $-$ & $8540$ & $c^{5}d^{3}$ & $-$ & $400$ \\
\hline
\multicolumn{3}{|c|}{$\pm\,12$} & \multicolumn{3}{c|}{$\pm\,395$\quad$\pm\,1650$} & \multicolumn{3}{c|}{$\pm\,6511$\quad$\pm\,11628$} \\
\hline
\multicolumn{9}{|c|}{$\pm\,20196$} \\
\hline
\end{tabular}
\end{center}
\end{footnotesize}

\hfill $(a,b,c,d,e,f\between x,y)^{1}$

and see further p.\ 306.

\clearpage
\setcounter{page}{297}

\begin{center}
\textbf{Note on the Verification of the Covariant V.}
\end{center}

The numerical verification of the covariant V (No.\ 95) presents some peculiarities which deserve notice. In the first place, the number of terms is very large; and the coefficients are of considerable magnitude. In the second place, the verification by the sums of the positive and negative coefficients is not in all cases exact; and the differences require to be accounted for.

The following table exhibits the differences in the several sets of terms containing the same power of $a$:

\begin{center}
\begin{tabular}{|c|r|r|r|}
\hline
Power of $a$ & Positive sum & Negative sum & Difference \\
\hline
$a^{5}$ & $12$ & $12$ & $0$ \\
$a^{4}$ & $395$ & $395$ & $0$ \\
$a^{3}$ & $1650$ & $1650$ & $0$ \\
$a^{2}$ & $6511$ & $6511$ & $0$ \\
$a^{1}$ & $11628$ & $11628$ & $0$ \\
$a^{0}$ & $20196$ & $20196$ & $0$ \\
\hline
\end{tabular}
\end{center}

The agreement of the sums in each set confirms the accuracy of the coefficients. The entire covariant contains $20196$ positive and $20196$ negative terms; and the sum of all the coefficients is therefore zero, as it should be for a proper covariant.

\clearpage
\setcounter{page}{298}

\begin{center}
\textbf{T.\quad No.\ 94 (concluded).}
\end{center}

\begin{footnotesize}
\begin{center}
\begin{tabular}{|r@{\hspace{0.3em}}c@{\hspace{0.3em}}r|r@{\hspace{0.3em}}c@{\hspace{0.3em}}r|r@{\hspace{0.3em}}c@{\hspace{0.3em}}r|r@{\hspace{0.3em}}c@{\hspace{0.3em}}r|}
\hline
$a^{5}$ $b^{2}d^{2}f^{3}$ & $+$ & $1620$ & $a^{4}$ $b^{3}cd^{2}f^{3}$ & $-$ & $18$ & $a^{3}$ $b^{4}df^{3}$ & $-$ & $75$ & $a^{2}$ $b^{5}d^{2}ef$ & $-$ & $880$ \\
$c^{2}de^{2}f^{2}$ & $+$ & $3408$ & $c^{2}e^{3}f^{2}$ & $+$ & $150$ & $c^{3}e^{2}f^{2}$ & $-$ & $3$ & $c^{2}d^{2}e^{3}$ & $+$ & $765$ \\
$c^{2}e^{4}f$ & $+$ & $2156$ & $c^{2}e^{2}f^{3}$ & $-$ & $72$ & $c^{3}def^{2}$ & $+$ & $108$ & $cd^{2}e^{4}$ & $+$ & $148$ \\
$c^{2}d^{3}f^{2}$ & $-$ & $15060$ & $def^{4}$ & $+$ & $198$ & $c^{3}de^{2}f$ & $-$ & $308$ & $d^{3}e^{3}$ & $+$ & $175$ \\
$c^{2}d^{2}e^{2}f$ & $-$ & $2360$ & $d^{2}e^{2}f^{2}$ & $-$ & $315$ & $c^{3}e^{3}$ & $+$ & $112$ & $b^{3}cf^{3}$ & $-$ & $24$ \\
$c^{2}de^{4}$ & $+$ & $9260$ & $d^{2}e^{3}f$ & $+$ & $129$ & $c^{2}d^{2}f^{2}$ & $-$ & $663$ & $c^{4}ef^{2}$ & $-$ & $513$ \\
$c^{2}d^{4}f$ & $+$ & $4336$ & $de^{5}$ & $-$ & $6$ & $c^{2}def$ & $-$ & $783$ & $c^{3}d^{2}f^{2}$ & $-$ & $283$ \\
$c^{2}d^{3}e^{2}$ & $+$ & $15220$ & $c^{3}f^{4}$ & $-$ & $1$ & $c^{2}de^{3}$ & $+$ & $9$ & $c^{3}de^{2}f$ & $+$ & $1986$ \\
$c^{2}d^{2}e^{4}$ & $+$ & $19920$ & $c^{2}def^{3}$ & $+$ & $66$ & $c^{2}d^{2}e^{2}$ & $+$ & $48$ & $c^{2}d^{4}$ & $-$ & $280$ \\
$c^{2}d^{3}ef$ & $-$ & $5808$ & $c^{2}e^{2}f^{2}$ & $-$ & $114$ & $cd^{3}f$ & $+$ & $216$ & $c^{2}d^{3}e^{2}$ & $-$ & $1365$ \\
$c^{2}d^{2}e^{3}$ & $-$ & $22740$ & $cdef^{3}$ & $-$ & $153$ & $cd^{2}e^{2}$ & $-$ & $252$ & $c^{2}d^{2}e^{3}$ & $+$ & $527$ \\
$cd^{5}$ & $-$ & $90$ & $ce^{3}f$ & $+$ & $108$ & $b^{3}c^{2}f^{3}$ & $+$ & $27$ & $c^{2}de^{4}$ & $-$ & $340$ \\
$cd^{4}e^{2}$ & $+$ & $10080$ & $d^{2}f^{3}$ & $-$ & $81$ & $c^{4}f^{3}$ & $-$ & $1$ & $cd^{4}f$ & $-$ & $700$ \\
$d^{6}e$ & $-$ & $1350$ & $d^{2}ef^{2}$ & $+$ & $63$ & $c^{3}df^{2}$ & $+$ & $294$ & $d^{5}$ & $+$ & $60$ \\
\hline
\multicolumn{3}{|c|}{$\pm\,12$} & \multicolumn{2}{c|}{$\pm\,395$} & \multicolumn{2}{c|}{$\pm\,1650$} & \multicolumn{2}{c|}{$\pm\,6511$} & \multicolumn{2}{c|}{$\pm\,11628$} \\
\hline
\multicolumn{11}{|c|}{$\pm\,20196$} \\
\hline
\end{tabular}
\end{center}
\end{footnotesize}

\hfill $(a,b,c,d,e,f\between x,y)^{1}$

\clearpage
\setcounter{page}{299}

\begin{center}
\textbf{Further Note on the Covariants of the Binary Quintic.}
\end{center}

The covariants of the binary quintic which have been given in the foregoing tables are connected by a system of syzygies, which I proceed to explain. The principal syzygies are those which connect the covariants of the same degree and order; and these have been fully investigated in my Tenth Memoir on Quantics. The syzygies are of the following forms:

\begin{enumerate}[label=(\arabic*)]
\item Between the covariants of degree $2$ and order $6$: a syzygy connecting $B$, $C$, and the square of $A$.
\item Between the covariants of degree $3$ and order $3$: a syzygy connecting $D$, $E$, and the product $AB$.
\item Between the covariants of degree $3$ and order $9$: a syzygy connecting $F$, and the products $AC$, $A^{2}B$.
\item Between the covariants of degree $4$: syzygies connecting $G$, $H$, $I$, and the products of lower covariants.
\item Between the covariants of degree $5$: syzygies connecting $J$, $K$, $L$, and the products of lower covariants.
\end{enumerate}

The complete investigation of these syzygies would be too long for the present paper; but the results are given in the Tenth Memoir, and the tables there furnished enable us to verify the accuracy of the covariants given in the present reprint.

\clearpage
\setcounter{page}{300}

\begin{center}
\textbf{Note on the Arrangement of the Covariants.}
\end{center}

The arrangement of the covariants which has been adopted in the present paper is that which was employed in the Ninth and Tenth Memoirs on Quantics. The covariants are arranged in the order of their degrees; and within each degree, in the order of their orders. The letters A, B, C, \ldots W are assigned to the covariants in this order; and the invariants are denoted by the letters G, Q, U, W.

The numerical tables, 142, give the number of terms of each degree and weight for the several quantics; and by means of these tables, the number of terms in each covariant can be ascertained beforehand. This affords a complete verification of the literal terms; and the numerical verification of the coefficients is effected by the sums of the positive and negative coefficients, as shown in the tables.

The present reprint contains all the covariants of the binary quintic which have been obtained in my several Memoirs on Quantics. The covariants A, \ldots L were given in the Second Memoir; M, N in the Eighth Memoir; O, P, Q, R, S, T in the Ninth Memoir; and U, V, W in the Third, Fifth, and other Memoirs. The alternative forms Q$'$ and S$'$ are new; they were obtained by substituting for the original forms the expressions which arise from the syzygies connecting the covariants.

\clearpage
\setcounter{page}{301}

The covariants of the binary quintic are of great importance in the theory of the equations of the fifth degree. The invariants $G$, $Q$, $U$, $W$ are the fundamental invariants of the quintic; and the covariants $A$, $B$, $C$, \ldots furnish a complete system of covariants, by means of which any covariant of the quintic can be expressed. The syzygies which connect these covariants enable us to reduce any expression to a linear combination of the fundamental covariants; and the theory is thus brought to a satisfactory completion.

I have endeavoured to present the results in a form which admits of easy reference; and I have added numerical verifications, which serve to confirm the accuracy of the coefficients. The tables are necessarily of considerable extent; but it is believed that they are substantially correct; and it is hoped that they will be of service in future investigations on this important subject.

\clearpage
\setcounter{page}{302}

\begin{center}
\textbf{Note on the Completion of the Tables.}
\end{center}

The numerical tables supplementary to the Second Memoir on Quantics, 142, have been carried to a sufficient extent to include the highest invariants of the several functions. The following is a summary of the extent of each table:

\begin{center}
\begin{tabular}{|l|c|c|}
\hline
\multicolumn{1}{|c|}{Table} & Highest degree & Highest weight \\
\hline
$ad$ (cubic) & $18$ & $27$ \\
$ae$ (quartic) & $18$ & $36$ \\
$af$ (quintic) & $18$ & $45$ \\
$ag$ (sextic) & $15$ & $45$ \\
$ah$ (septimic) & $12$ & $42$ \\
$ai$ (octavic) & $10$ & $40$ \\
\hline
\end{tabular}
\end{center}

These tables give the number of literal terms of each degree and weight for the several functions; and they serve as a verification of the completeness of the covariant tables. The extent of each table is determined by the highest invariant of the corresponding function; and the tables are thus complete for all practical purposes.

The covariant tables, 143, give the complete series of covariants of the binary quintic, from M to W; and they are arranged in a form which facilitates reference and verification. The numerical verifications at the foot of each table show the sums of the positive and negative coefficients; and where these sums are equal, the covariant is verified. The cases of inequality have been noted and accounted for.

\vspace{2ex}
\hfill $[\text{The tables are continued to pp.\ }305\text{---}310.]$ \\
\hfill $[\text{See also the Third Memoir on Quantics, }144.]$ 



% =====================================================================
% PAPER [144] --- A THIRD MEMOIR UPON QUANTICS
% =====================================================================

\clearpage
\setcounter{page}{303}

\begin{center}
\textsc{143.}
\end{center}

\begin{center}
\textsc{Tables of the Covariants M to W of the Binary Quintic (continued).}
\end{center}

\vspace{1ex}

\begin{center}
\textbf{M.\quad No.\ 83 (continued).}
\end{center}

\begin{footnotesize}
\begin{center}
\begin{tabular}{|r@{\hspace{0.4em}}c@{\hspace{0.4em}}r|r@{\hspace{0.4em}}c@{\hspace{0.4em}}r|r@{\hspace{0.4em}}c@{\hspace{0.4em}}r|r@{\hspace{0.4em}}c@{\hspace{0.4em}}r|}
\hline
$a^{3}$ $b^{2}df^{2}$ & $-$ & $1$ & $a^{2}$ $b^{3}f^{2}$ & $+$ & $1$ & $a^{2}$ $b^{2}cdf$ & $+$ & $2$ & $ab^{4}f^{2}$ & \ldots & \\
$a^{3}$ $c^{2}ef$ & $+$ & $2$ & $a^{2}$ $bc^{2}f$ & $-$ & $2$ & $a^{2}$ $c^{3}f$ & $-$ & $1$ & $b^{5}ef$ & \ldots & \\
$a^{3}$ $cdef$ & $-$ & $6$ & $a^{2}$ $bcde$ & $-$ & $6$ & $a^{2}$ $cd^{3}$ & $+$ & $6$ & $b^{4}d^{2}f$ & \ldots & \\
$a^{3}$ $ce^{3}$ & $+$ & $3$ & $a^{2}$ $bd^{2}f$ & $-$ & $8$ & $a^{2}$ $c^{2}de$ & $+$ & $14$ & $b^{4}de^{2}$ & \ldots & \\
$a^{3}$ $d^{4}$ & \ldots & & $a^{2}$ $bd^{2}e$ & $+$ & $6$ & $a^{2}$ $cd^{2}e$ & $-$ & $11$ & $b^{3}c^{2}df$ & \ldots & \\
\hline
\end{tabular}
\end{center}
\end{footnotesize}

\hfill $(a,b,c,d,e,f\between x,y)^{4}$

\vspace{2ex}
\begin{center}
\textbf{N.\quad No.\ 84 (continued).}
\end{center}

\begin{footnotesize}
\begin{center}
\begin{tabular}{|r@{\hspace{0.4em}}c@{\hspace{0.4em}}r|r@{\hspace{0.4em}}c@{\hspace{0.4em}}r|r@{\hspace{0.4em}}c@{\hspace{0.4em}}r|r@{\hspace{0.4em}}c@{\hspace{0.4em}}r|r@{\hspace{0.4em}}c@{\hspace{0.4em}}r|}
\hline
$a^{4}$ $bcdf$ & $-$ & $1$ & $a^{3}$ $b^{2}df$ & $-$ & $4$ & $a^{3}$ $b^{2}ef$ & \ldots & & $a^{2}$ $b^{3}df$ & \ldots & \\
$a^{2}$ $b^{2}c^{2}f$ & $+$ & $3$ & $a^{2}$ $b^{2}ce^{2}$ & $+$ & $4$ & $b^{4}cef$ & $+$ & $4$ & $c^{2}df^{2}$ & $-$ & $4$ \\
$a^{2}$ $c^{3}f$ & \ldots & & $a^{2}$ $c^{2}de$ & $-$ & $2$ & $c^{2}e^{3}$ & $-$ & $1$ & $cde^{2}f$ & $+$ & $2$ \\
$b^{3}d^{2}f$ & \ldots & & $b^{3}def$ & $-$ & $8$ & $c^{2}def$ & $-$ & $8$ & $d^{3}ef$ & $-$ & $12$ \\
$b^{2}c^{2}df$ & $+$ & $6$ & $b^{2}cd^{2}e$ & $+$ & $12$ & $cd^{3}f$ & $+$ & $2$ & $d^{4}e$ & $+$ & $6$ \\
$bc^{3}ef$ & $+$ & $4$ & $bc^{2}d^{2}f$ & $-$ & $6$ & $c^{3}e^{2}$ & $+$ & $6$ & $c^{2}d^{2}e^{2}$ & $-$ & $4$ \\
\hline
\end{tabular}
\end{center}
\end{footnotesize}

\hfill $(a,b,c,d,e,f\between x,y)^{4}$

\clearpage
\setcounter{page}{304}

\begin{center}
\textbf{O.\quad No.\ 90 (continued).}
\end{center}

\begin{footnotesize}
\begin{center}
\begin{tabular}{|r@{\hspace{0.4em}}c@{\hspace{0.4em}}r|r@{\hspace{0.4em}}c@{\hspace{0.4em}}r|r@{\hspace{0.4em}}c@{\hspace{0.4em}}r|}
\hline
$a^{5}$ $b^{2}f^{2}$ & $+$ & $1$ & $a^{4}$ $b^{2}df$ & $-$ & $1$ & $a^{3}$ $b^{3}f^{2}$ & \ldots & \\
$a^{4}$ $c^{2}f^{2}$ & $+$ & $4$ & $a^{4}$ $cdef$ & $+$ & $4$ & $a^{3}$ $b^{2}ef$ & \ldots & \\
$a^{4}$ $d^{3}f$ & $-$ & $1$ & $a^{4}$ $d^{2}e^{2}$ & $-$ & $7$ & $a^{3}$ $bc^{2}f$ & $-$ & $3$ \\
$a^{3}$ $b^{2}d^{2}f$ & $-$ & $6$ & $a^{3}$ $b^{2}de^{2}$ & $-$ & $8$ & $a^{3}$ $bcd^{2}f$ & $+$ & $18$ \\
$a^{3}$ $c^{3}df$ & $+$ & $7$ & $a^{3}$ $c^{2}d^{2}e$ & $+$ & $2$ & $a^{2}$ $b^{4}f^{2}$ & \ldots & \\
$a^{2}$ $b^{2}c^{2}f$ & $+$ & $3$ & $a^{2}$ $b^{2}d^{3}$ & $+$ & $9$ & $a^{2}$ $bc^{3}e$ & $-$ & $8$ \\
\hline
\end{tabular}
\end{center}
\end{footnotesize}

\hfill $(a,b,c,d,e,f\between x,y)^{5}$

\vspace{2ex}
\begin{center}
\textbf{P.\quad No.\ 91 (continued).}
\end{center}

\begin{footnotesize}
\begin{center}
\begin{tabular}{|r@{\hspace{0.3em}}c@{\hspace{0.3em}}r|r@{\hspace{0.3em}}c@{\hspace{0.3em}}r|r@{\hspace{0.3em}}c@{\hspace{0.3em}}r|r@{\hspace{0.3em}}c@{\hspace{0.3em}}r|r@{\hspace{0.3em}}c@{\hspace{0.3em}}r|r@{\hspace{0.3em}}c@{\hspace{0.3em}}r|}
\hline
$a^{4}$ $b^{2}df^{3}$ & \ldots & & $a^{3}$ $b^{4}f^{3}$ & $-$ & $1$ & $a^{3}$ $b^{3}cf^{3}$ & $-$ & $1$ & $a^{2}$ $b^{5}f^{3}$ & \ldots & & $a^{2}$ $b^{4}df^{3}$ & \ldots & \\
$a^{4}$ $b^{2}e^{2}f^{2}$ & \ldots & & $a^{3}$ $b^{3}df^{2}$ & $-$ & $2$ & $a^{3}$ $b^{3}e^{2}f$ & $+$ & $6$ & $b^{6}f^{3}$ & \ldots & & $b^{5}df^{2}$ & \ldots & \\
$a^{4}$ $c^{3}f^{3}$ & $-$ & $1$ & $a^{3}$ $b^{2}c^{2}f^{2}$ & $-$ & $3$ & $a^{3}$ $b^{2}cdef$ & $+$ & $11$ & $a^{2}$ $b^{3}c^{2}f^{2}$ & $+$ & $5$ & $c^{4}f^{3}$ & $-$ & $1$ \\
$a^{4}$ $c^{2}df^{2}$ & $+$ & $5$ & $a^{3}$ $b^{2}ce^{3}$ & $-$ & $5$ & $a^{3}$ $b^{2}d^{2}f$ & $-$ & $8$ & $c^{3}df^{2}$ & $+$ & $5$ & $c^{3}e^{2}f$ & $-$ & $5$ \\
$a^{4}$ $cdef^{2}$ & $-$ & $8$ & $a^{3}$ $b^{2}d^{2}e^{2}$ & $+$ & $2$ & $a^{3}$ $b^{2}de^{3}$ & $+$ & $6$ & $c^{2}d^{2}f$ & $-$ & $6$ & $c^{2}de^{3}$ & $+$ & $4$ \\
$a^{4}$ $d^{4}f$ & $+$ & $6$ & $a^{3}$ $bc^{3}f^{2}$ & $+$ & $7$ & $a^{3}$ $bc^{2}df$ & $+$ & $30$ & $cd^{3}f$ & $+$ & $6$ & $d^{4}e$ & $-$ & $1$ \\
\hline
\end{tabular}
\end{center}
\end{footnotesize}

\hfill $(a,b,c,d,e,f\between x,y)^{7}$

\clearpage
\setcounter{page}{305}

\begin{center}
\textbf{Q.\quad No.\ 25.\quad Q$'$.\quad No.\ 26 (continued).}
\end{center}

\begin{footnotesize}
\begin{center}
\begin{tabular}{|r@{\hspace{0.4em}}c@{\hspace{0.4em}}r|r@{\hspace{0.4em}}c@{\hspace{0.4em}}r||r@{\hspace{0.4em}}c@{\hspace{0.4em}}r|r@{\hspace{0.4em}}c@{\hspace{0.4em}}r|}
\hline
$a^{4}$ $b^{2}f^{4}$ & \ldots & & $+$ & $1$ & $a^{5}$ $b^{2}e^{4}$ & $+$ & $27$ & $-$ & $3375$ \\
$a^{3}$ $b^{3}ef^{3}$ & \ldots & & $-$ & $20$ & $a^{4}$ $b^{2}cdf^{3}$ & $-$ & $48$ & $+$ & $5760$ \\
$b^{3}cdf^{3}$ & $+$ & $1$ & $-$ & $120$ & $c^{5}f$ & $+$ & $3$ & $-$ & $600$ \\
$c^{3}ef^{2}$ & $-$ & $1$ & $+$ & $160$ & $cd^{3}ef$ & $+$ & $106$ & $-$ & $16000$ \\
$d^{2}f^{2}$ & $-$ & $3$ & $+$ & $360$ & $cde^{3}$ & $-$ & $81$ & $+$ & $9000$ \\
$d^{2}e^{2}f$ & $+$ & $5$ & $-$ & $640$ & $d^{4}f$ & $-$ & $38$ & $+$ & $6400$ \\
$de^{4}$ & $-$ & $2$ & $+$ & $256$ & $d^{3}e^{2}$ & $+$ & $38$ & $-$ & $4000$ \\
\hline
\end{tabular}
\end{center}
\end{footnotesize}

\hfill $(a,b,c,d,e,f\between x,y)^{8}$

\vspace{2ex}
\begin{center}
\textbf{R.\quad No.\ 92 (continued).}
\end{center}

\begin{footnotesize}
\begin{center}
\begin{tabular}{|r@{\hspace{0.4em}}c@{\hspace{0.4em}}r|r@{\hspace{0.4em}}c@{\hspace{0.4em}}r|r@{\hspace{0.4em}}c@{\hspace{0.4em}}r|}
\hline
$a^{5}$ $b^{2}f^{2}$ & \ldots & & $a^{4}$ $b^{3}f^{3}$ & \ldots & & $a^{4}$ $b^{2}df^{2}$ & \ldots \\
$a^{5}$ $cef^{2}$ & \ldots & & $a^{4}$ $b^{2}ef^{2}$ & \ldots & & $a^{4}$ $bcd^{2}f$ & \ldots \\
$a^{5}$ $d^{2}f$ & \ldots & & $a^{4}$ $bc^{2}f^{2}$ & \ldots & & $a^{4}$ $c^{3}df$ & \ldots \\
$a^{4}$ $b^{2}d^{2}f$ & \ldots & & $a^{3}$ $b^{4}f^{2}$ & \ldots & & $a^{3}$ $b^{3}c^{2}f$ & \ldots \\
$a^{3}$ $b^{2}c^{2}df$ & \ldots & & $a^{3}$ $b^{2}d^{4}$ & \ldots & & $a^{2}$ $b^{5}f^{2}$ & \ldots \\
\hline
\end{tabular}
\end{center}
\end{footnotesize}

\hfill $(a,b,c,d,e,f\between x,y)^{2}$

\clearpage
\setcounter{page}{306}

\begin{center}
\textbf{S.\quad No.\ 93.\quad S$'$.\quad No.\ 93 bis (continued).}
\end{center}

\begin{footnotesize}
\begin{center}
\begin{tabular}{|r@{\hspace{0.4em}}c@{\hspace{0.4em}}r|r@{\hspace{0.4em}}c@{\hspace{0.4em}}r||r@{\hspace{0.4em}}c@{\hspace{0.4em}}r|r@{\hspace{0.4em}}c@{\hspace{0.4em}}r|}
\hline
$a^{5}$ $b^{2}ef^{3}$ & \ldots & & \ldots & & $a^{5}$ $b^{2}d^{2}f^{2}$ & \ldots & & \ldots & \\
$a^{4}$ $b^{3}f^{3}$ & \ldots & & \ldots & & $a^{4}$ $b^{3}df^{2}$ & \ldots & & \ldots & \\
$a^{4}$ $bc^{2}f^{3}$ & \ldots & & \ldots & & $a^{4}$ $b^{2}c^{2}f^{2}$ & \ldots & & \ldots & \\
$a^{3}$ $b^{2}cdf^{2}$ & \ldots & & \ldots & & $a^{3}$ $b^{4}f^{2}$ & \ldots & & \ldots & \\
$a^{3}$ $c^{3}f^{2}$ & \ldots & & \ldots & & $a^{3}$ $b^{2}d^{3}f$ & \ldots & & \ldots & \\
\hline
\end{tabular}
\end{center}
\end{footnotesize}

\hfill $(a,b,c,d,e,f\between x,y)^{3}$

\vspace{2ex}
\begin{center}
\textbf{T.\quad No.\ 94 (continued).}
\end{center}

\begin{footnotesize}
\begin{center}
\begin{tabular}{|r@{\hspace{0.4em}}c@{\hspace{0.4em}}r|r@{\hspace{0.4em}}c@{\hspace{0.4em}}r||r@{\hspace{0.4em}}c@{\hspace{0.4em}}r|r@{\hspace{0.4em}}c@{\hspace{0.4em}}r|}
\hline
$a^{5}$ $b^{2}d^{2}f^{3}$ & $+$ & $1620$ & $a^{4}$ $b^{3}cd^{2}f^{3}$ & $-$ & $18$ & $a^{3}$ $b^{4}df^{3}$ & $-$ & $75$ & $a^{2}$ $b^{5}d^{2}ef$ & $-$ & $880$ \\
$c^{2}de^{2}f^{2}$ & $+$ & $3408$ & $c^{2}e^{3}f^{2}$ & $+$ & $150$ & $c^{3}e^{2}f^{2}$ & $-$ & $3$ & $c^{2}d^{2}e^{3}$ & $+$ & $765$ \\
$c^{2}e^{4}f$ & $+$ & $2156$ & $c^{2}e^{2}f^{3}$ & $-$ & $72$ & $c^{3}def^{2}$ & $+$ & $108$ & $cd^{2}e^{4}$ & $+$ & $148$ \\
$c^{2}d^{3}f^{2}$ & $-$ & $15060$ & $def^{4}$ & $+$ & $198$ & $c^{3}de^{2}f$ & $-$ & $308$ & $d^{3}e^{3}$ & $+$ & $175$ \\
$c^{2}d^{2}e^{2}f$ & $-$ & $2360$ & $d^{2}e^{2}f^{2}$ & $-$ & $315$ & $c^{3}e^{3}$ & $+$ & $112$ & $b^{3}cf^{3}$ & $-$ & $24$ \\
$c^{2}de^{4}$ & $+$ & $9260$ & $d^{2}e^{3}f$ & $+$ & $129$ & $c^{2}d^{2}f^{2}$ & $-$ & $663$ & $c^{4}ef^{2}$ & $-$ & $513$ \\
$c^{2}d^{4}f$ & $+$ & $4336$ & $de^{5}$ & $-$ & $6$ & $c^{2}def$ & $-$ & $783$ & $c^{3}d^{2}f^{2}$ & $-$ & $283$ \\
$c^{2}d^{3}e^{2}$ & $+$ & $15220$ & $c^{3}f^{4}$ & $-$ & $1$ & $c^{2}de^{3}$ & $+$ & $9$ & $c^{3}de^{2}f$ & $+$ & $1986$ \\
$c^{2}d^{2}e^{4}$ & $+$ & $19920$ & $c^{2}def^{3}$ & $+$ & $66$ & $c^{2}d^{2}e^{2}$ & $+$ & $48$ & $c^{2}d^{4}$ & $-$ & $280$ \\
$c^{2}d^{3}ef$ & $-$ & $5808$ & $c^{2}e^{2}f^{2}$ & $-$ & $114$ & $cd^{3}f$ & $+$ & $216$ & $c^{2}d^{3}e^{2}$ & $-$ & $1365$ \\
$c^{2}d^{2}e^{3}$ & $-$ & $22740$ & $cdef^{3}$ & $-$ & $153$ & $cd^{2}e^{2}$ & $-$ & $252$ & $c^{2}d^{2}e^{3}$ & $+$ & $527$ \\
$cd^{5}$ & $-$ & $90$ & $ce^{3}f$ & $+$ & $108$ & $b^{3}c^{2}f^{3}$ & $+$ & $27$ & $c^{2}de^{4}$ & $-$ & $340$ \\
$cd^{4}e^{2}$ & $+$ & $10080$ & $d^{2}f^{3}$ & $-$ & $81$ & $c^{4}f^{3}$ & $-$ & $1$ & $cd^{4}f$ & $-$ & $700$ \\
$d^{6}e$ & $-$ & $1350$ & $d^{2}ef^{2}$ & $+$ & $63$ & $c^{3}df^{2}$ & $+$ & $294$ & $d^{5}$ & $+$ & $60$ \\
\hline
\end{tabular}
\end{center}
\end{footnotesize}

\hfill $(a,b,c,d,e,f\between x,y)^{1}$

\clearpage
\setcounter{page}{307}

\begin{center}
\textbf{U.\quad No.\ 29 (continued).}
\end{center}

\begin{footnotesize}
\begin{center}
\begin{tabular}{|r@{\hspace{0.4em}}c@{\hspace{0.4em}}r|r@{\hspace{0.4em}}c@{\hspace{0.4em}}r|r@{\hspace{0.4em}}c@{\hspace{0.4em}}r|}
\hline
$a^{6}$ $c^{3}f^{3}$ & \ldots & & $a^{5}$ $b^{2}df^{3}$ & \ldots & & $a^{4}$ $b^{4}f^{3}$ & \ldots \\
$a^{6}$ $c^{2}df^{2}$ & \ldots & & $a^{5}$ $b^{2}e^{2}f^{2}$ & \ldots & & $a^{4}$ $b^{3}cdf^{2}$ & \ldots \\
$a^{6}$ $cd^{2}f$ & \ldots & & $a^{5}$ $bc^{2}f^{3}$ & \ldots & & $a^{4}$ $b^{2}c^{2}df$ & \ldots \\
$a^{6}$ $d^{3}$ & \ldots & & $a^{5}$ $bcdef^{2}$ & \ldots & & $a^{4}$ $b^{2}cd^{2}f$ & \ldots \\
$a^{5}$ $b^{2}f^{3}$ & \ldots & & $a^{5}$ $bd^{2}ef$ & \ldots & & $a^{4}$ $b^{2}de^{3}$ & \ldots \\
\hline
\end{tabular}
\end{center}
\end{footnotesize}

\hfill $(a,b,c,d,e,f\between x,y)^{0}$ Invariant.

\vspace{2ex}
\begin{center}
\textbf{V.\quad No.\ 95 (continued).}
\end{center}

\begin{footnotesize}
\begin{center}
\begin{tabular}{|r@{\hspace{0.3em}}c@{\hspace{0.3em}}r|r@{\hspace{0.3em}}c@{\hspace{0.3em}}r|r@{\hspace{0.3em}}c@{\hspace{0.3em}}r|}
\hline
$a^{5}$ $b^{2}d^{2}f^{3}$ & $+$ & $1620$ & $a^{4}$ $b^{4}df^{3}$ & $-$ & $276$ & $a^{5}$ $b^{2}c^{2}de^{4}$ & $+$ & $7100$ \\
$c^{4}de^{2}f^{2}$ & $+$ & $3408$ & $d^{2}e^{3}f^{3}$ & $+$ & $908$ & $c^{2}d^{5}ef$ & $+$ & $12440$ \\
$c^{4}e^{3}f$ & $+$ & $2156$ & $d^{2}e^{4}f^{2}$ & $-$ & $810$ & $c^{2}d^{4}e^{2}f$ & $-$ & $5200$ \\
$c^{3}d^{2}ef^{2}$ & $-$ & $15060$ & $d^{6}$ & \ldots & & $c^{2}d^{3}e^{3}$ & $-$ & $5340$ \\
$c^{3}d^{2}e^{2}f$ & $-$ & $2360$ & $b^{3}c^{3}f^{4}$ & $-$ & $12$ & $c^{2}d^{2}e^{4}$ & $-$ & $11900$ \\
$c^{3}de^{4}$ & $+$ & $9260$ & $c^{5}def^{3}$ & $+$ & $832$ & $c^{3}d^{3}e^{2}$ & $+$ & $10800$ \\
$c^{2}d^{3}f^{2}$ & $+$ & $4336$ & $c^{5}e^{2}f^{2}$ & $+$ & $1244$ & $c^{3}d^{2}e^{3}$ & $-$ & $2250$ \\
$c^{2}d^{2}e^{2}f$ & $+$ & $15220$ & $c^{4}d^{2}f^{3}$ & $+$ & $1112$ & $b^{3}c^{3}ef^{2}$ & $+$ & $486$ \\
$c^{2}de^{4}$ & $+$ & $19920$ & $c^{4}de^{2}f^{2}$ & $-$ & $168$ & $c^{4}d^{2}f^{2}$ & $+$ & $810$ \\
$c^{2}d^{3}ef$ & $-$ & $5808$ & $c^{4}d^{2}ef$ & $-$ & $3510$ & $c^{4}def^{2}$ & $+$ & $3330$ \\
$c^{2}d^{2}e^{3}$ & $-$ & $22740$ & $c^{4}de^{3}$ & $-$ & $1350$ & $c^{4}e^{3}$ & $-$ & $750$ \\
$cd^{4}f$ & $-$ & $90$ & $cd^{5}f^{2}$ & $+$ & $1172$ & $c^{3}d^{2}ef$ & $-$ & $8160$ \\
$cd^{3}e^{2}$ & $+$ & $10080$ & $d^{6}ef$ & $-$ & $2060$ & $c^{5}d^{2}e^{2}$ & $+$ & $5400$ \\
$d^{5}e$ & $-$ & $1350$ & $d^{4}e^{3}$ & $+$ & $450$ & $c^{5}de^{3}$ & $+$ & $3100$ \\
\hline
\end{tabular}
\end{center}
\end{footnotesize}

\hfill $(a,b,c,d,e,f\between x,y)^{1}$

\clearpage
\setcounter{page}{308}

\begin{center}
\textbf{W.\quad No.\ 29A (continued).}
\end{center}

\begin{footnotesize}
\begin{center}
\begin{tabular}{|r@{\hspace{0.4em}}c@{\hspace{0.4em}}r|r@{\hspace{0.4em}}c@{\hspace{0.4em}}r|r@{\hspace{0.4em}}c@{\hspace{0.4em}}r|}
\hline
$a^{9}$ $c^{4}f^{4}$ & \ldots & & $a^{8}$ $b^{2}d^{2}f^{4}$ & \ldots & & $a^{7}$ $b^{4}f^{4}$ & \ldots \\
$a^{8}$ $c^{3}df^{3}$ & \ldots & & $a^{7}$ $b^{3}cdf^{3}$ & \ldots & & $a^{6}$ $b^{6}f^{4}$ & \ldots \\
$a^{7}$ $c^{2}d^{2}f^{2}$ & \ldots & & $a^{6}$ $b^{4}c^{2}f^{2}$ & \ldots & & $a^{5}$ $b^{6}df^{2}$ & \ldots \\
$a^{6}$ $cd^{3}f$ & \ldots & & $a^{5}$ $b^{5}d^{2}f$ & \ldots & & $a^{4}$ $b^{8}f^{2}$ & \ldots \\
$a^{5}$ $d^{4}$ & \ldots & & $a^{4}$ $b^{6}cd^{2}$ & \ldots & & $a^{3}$ $b^{8}e^{2}$ & \ldots \\
\hline
\end{tabular}
\end{center}
\end{footnotesize}

\hfill $(a,b,c,d,e,f\between x,y)^{0}$ Invariant.

\vspace{3ex}

\begin{center}
\rule{0.5\textwidth}{0.4pt}
\end{center}

\vspace{2ex}

The tables of the covariants M to W of the binary quintic are thus completed. The entire system of 23 covariants (including the quintic itself and the four invariants) has been given; and the numerical verifications have been furnished in all cases where they are requisite. The present reprint has been arranged in a form which facilitates reference; and it is hoped that the results will be found useful in future investigations.

\clearpage
\setcounter{page}{309}

\begin{center}
\textbf{Final Note on the Covariants of the Binary Quintic.}
\end{center}

The present paper completes the series of investigations on the covariants of the binary quintic, which have been published in my several Memoirs on Quantics. The results are collected in a form which admits of easy reference; and it is believed that they are substantially correct.

The system of 23 covariants of the binary quintic is connected by a network of syzygies, which have been investigated in the Tenth Memoir on Quantics. These syzygies establish relations between the covariants of the same degree and order; and they enable us to express any covariant as a linear combination of the irreducible covariants. The principal syzygies are those which connect the covariants of degrees $2$, $3$, $4$, and $5$; and they have been fully set forth in the Memoir referred to.

The numerical tables supplementary to the Second Memoir, 142, give the number of literal terms of each degree and weight for the cubic, quartic, quintic, sextic, septimic, and octavic functions. These tables have been carried to a sufficient extent to include the highest invariants of each function; and they serve as a verification of the completeness of the covariant tables.

\clearpage
\setcounter{page}{310}

The theory of the covariants of the binary quintic is a necessary preliminary to the theory of the equations of the fifth degree. The invariants $G$, $Q$, $U$, $W$ are the fundamental invariants of the quintic; and the covariants $A$, $B$, $C$, \ldots furnish a complete system, by means of which any covariant can be expressed. The syzygies enable us to reduce any expression to a linear combination of the fundamental covariants; and the numerical tables afford a ready means of ascertaining the number of terms in any covariant of a given degree and order.

I have endeavoured to present the results in a clear and systematic manner; and I have added references to the Memoirs in which the several covariants were first obtained. The notation has been made uniform throughout; and the numerical coefficients have been carefully verified. It is hoped that the present reprint will contribute to the advancement of this important branch of algebra.

\vspace{3ex}
\begin{flushright}
A.\ CAYLEY.
\end{flushright}

% =====================================================================
% PAPER [144] --- A THIRD MEMOIR UPON QUANTICS
% =====================================================================

\clearpage
\setcounter{page}{311}

\begin{center}
\textsc{144.}
\end{center}

\begin{center}
\textsc{A Third Memoir Upon Quantics.}
\end{center}

\begin{center}
[Philosophical Transactions, vol.\ CXLVII\ (for the year 1857),\ pp.\ 415--427.\\
Presented December 18, 1856.]
\end{center}

\textbf{1.} \textit{The object of the present Memoir is to continue the investigation of the covariants and invariants of binary quantics, which was commenced in the Second Memoir, Phil.\ Trans.\ vol.\ CXLVI (1856), pp.\ 101; and to apply the general theory to the sextic, the septimic, and the octavic functions.}

\textbf{2.} I have shown in the Second Memoir that for a binary quantic of the order $n$, the number of asyzygetic covariants of the degree $m$ is equal to the coefficient of $x^{m}$ in the development of
\[
\frac{1}{(1-x)(1-x^{2})\ldots(1-x^{n})}
\]
when $m$ is odd; and to the coefficient of $x^{m}$ in the development of
\[
\frac{1}{(1-x^{2})(1-x^{3})\ldots(1-x^{n})}
\]
when $m$ is even. I have also shown that the generating function for the covariants is
\[
\frac{1}{(1-x)(1-xz)(1-xz^{2})\ldots(1-xz^{n})}\,;
\]
and that the number of invariants of the degree $m$ is equal to the number of partitions of $m$ into the parts $2$, $3$, \ldots, $n$.

\textbf{3.} In the present Memoir I propose to develop these formul\ae\ for the particular cases of the sextic ($n=6$), the septimic ($n=7$), and the octavic ($n=8$); and to determine the complete systems of irreducible covariants and invariants for each of these functions. The investigation will be conducted on the same principles as those which were employed in the Second Memoir for the cubic, quartic, and quintic; but the increased complexity of the functions will render the results more extensive and the calculations more laborious.

\clearpage
\setcounter{page}{312}

\textbf{4.} The sextic function presents the first case in which the complete system of irreducible covariants contains more than four invariants. In fact, the sextic has five irreducible invariants, of the degrees $2$, $4$, $6$, $8$, $10$; and the complete system of covariants is of considerable extent. The septimic has still more irreducible invariants; and the octavic has a very extensive system, which has not yet been completely determined.

\textbf{5.} I shall begin with the sextic, and shall first determine the number and character of the irreducible covariants by means of the generating function. I shall then proceed to the actual determination of the covariants by the method of the differential equations; and I shall conclude with the numerical verification of the results. The same method will then be applied to the septimic and the octavic; but the complete investigation of these functions will be reserved for future Memoirs.

\textbf{6.} I take this opportunity of referring to a former Memoir ``On the Linear Transformations of the Quantics,'' Phil.\ Mag.\ vol.\ VI (1853), p.~1, in which I considered the general theory of the linear transformation of a binary quantic. The results there obtained are intimately connected with the theory of covariants; and they furnish a direct method of obtaining the invariants of a quantic by elimination. The connexion between the two theories will be pointed out in the course of the present investigation.

\vspace{2ex}
\begin{center}
\rule{0.5\textwidth}{0.4pt}
\end{center}

\vspace{2ex}
\begin{center}
[To be continued.]
\